% 来源: 19c9fbc3-c752-8f62-8000-00009c915a0b_3.2.1_交换机设计.pdf
% 提取时间: 2026-02-28T22:51:46+08:00

3.2.1 交换机设计
1 交换机
在电话通信中，一次电话连接通常持续数秒甚至数分钟。然而在互联网交换机中，每次连接仅持续一个
数据包的传输时间。对于一个40字节的数据包，在40Gbps的速率下，这个时间是8纳秒。因此，为了以
线速运行，交换系统必须在最短的数据包到达时间内决定（即计算）应该匹配哪些输入和输出链路。这
使得互联网交换机的控制部分（用于建立连接）比电话交换机更难构建。为此，
在数据面，大多数路由器在将可变大小的数据包发送到交换结构之前，会在内部将其分割为固定大小的
信元。
在控制面，根据数据到达情况，交换组件输入输出端口建立连接，可归结为解决一个二分匹配问题：路
由器必须在固定的信元到达时间（时隙）内，尽可能地将输入链路与输出链路进行匹配
1.1 发展历程：从总线到Crossbar




图1 网络交换机的演进，从带有共享CPU的共享总线交换机到每个线路卡带有专用转发引擎的CrossBar
交换机

                                                      1
在介绍基于总线和CrossBar（crossbar）的交换机之前，先了解一种基于共享内存的简单交换机实现方
式是有益的。这种设计的基本原理是：数据包从输入链路进入后被存储到共享内存中，随后从内存读取
并发送到对应的输出链路。这种技术在电话通信中的时隙交换中已有多年应用，同时也适用于小型网络
交换机。
然而，这种设计的主要瓶颈在于内存带宽。假设芯片有8个输入链路和8个输出链路，那么每个数据包需
要经历一次写入和一次读取操作，这意味着内存的运行速度必须达到单链路速度的16倍。为缓解这一问
题，可以通过增加内存访问宽度来优化。例如，输入链路上的串行比特流可以累积到一个移位寄存器
中，当完整的信元累积完成后，将其一次性写入宽字内存。类似地，数据包在读取时也会被加载到输出
移位寄存器中，随后逐位移出到输出链路。
尽管这种共享内存的设计易于在芯片上实现，并可通过流量控制和外部缓冲区扩展，但其扩展性有限。
主要原因是它无法处理超过信元大小的数据包。如果交换机同时收到多个发往不同目的地的小型数据
包，虽然可以将它们打包成一个宽字存储，但无法保证它们能同时被读出。
随着技术发展，交换机已经从简单的共享介质（如总线或共享内存）架构演进到更现代的CrossBar架
构。在早期的共享总线设计中，所有线路卡通过一条高速总线连接，每次只能允许一对线路卡通信。例
如，当线路卡1向线路卡2发送数据包时，其他线路卡之间的通信会被阻塞。此外，在早期路由器中，转
发决策由一个通用CPU负责，这虽然简化了软件设计，但由于通用CPU的低效指令执行和缺乏实时约束
能力，导致性能受限。同时，数据包需要两次经过总线（一次到CPU，一次从CPU到目标），进一步增
加了延迟。
为了缓解CPU瓶颈，后续设计引入了多CPU架构，每个CPU负责一部分线路卡的转发任务。这种方法提
高了吞吐量并降低了单个CPU的性能要求，但也可能引发数据包乱序问题。
尽管如此，共享总线仍然是系统的主要瓶颈。由于总线需要支持多个数据源和目的地，电气负载、信号
上升时间和反射等问题限制了其传输速度。为解决这些问题，CrossBar交换机应运而生。CrossBar交换
机采用一组2N条并行总线（每条源线路卡和目的线路卡各对应一条总线），形成矩阵结构。理论上，这
种设计比单总线架构快N倍，因为在最佳情况下，所有N条总线可以同时传输数据。然而，要充分发挥其
性能潜力，需要解决调度问题，即在每个时隙找到N对不冲突的源-目通信对。这一问题构成了本章研究
的核心内容。
1.2 CrossBar交换结构




                                                    2
图3 基本CrossBar交换机
最简单的CrossBar结构是一个由N条输入总线和N条输出总线组成的阵列，如图3所示。因此，如果线路
卡R希望向线路卡S发送数据，则输入总线R必须连接到输出总线S。实现这种连接的最简单方法是通过一
个“传输”晶体管，如图4所示。对于每一对输入和输出总线，如R和S，都有一个晶体管，当晶体管导通
时，就连接了这两条总线。这样的连接被称为交叉点。请注意，一个具有N个输入和N个输出的CrossBar
交换机有N²个交叉点，每个交叉点都需要一条来自调度器的控制线来控制其导通或关闭。




图4 通过导通连接两条总线的传输晶体管，将来自线路卡R的输入连接到线路卡S。现代CrossBar交换机
用多路复用器树取代了这种简单设计，以降低电容
虽然N²个交叉点看起来数量很多，但通过晶体管进行的超大规模集成电路（VLSI）实现使得引脚数量、
卡连接器技术等成为构建大型交换机时更具限制性的因素。因此，在2002年之前构建的大多数路由器和


                                                  3
交换机使用简单的CrossBar交换背板来支持16 - 32个端口。请注意，通过将输入总线R连接到所有希望
从R接收数据的输出总线，就可以轻松实现多播。然而，多播调度却很棘手。
在实际应用中，只有较旧的CrossBar设计才使用传输晶体管。这是因为随着端口数量的增加，总电容
（第2章）会变得非常大。这反过来又增加了信号传输延迟，在更高速度下这成为了一个问题。现代实现
方式通常在每个输出端使用大型多路复用器树或三态缓冲器。更高性能的系统甚至在交叉点处使用一些
内存（即一个门）对通过CrossBar结构的数据进行流水线处理。

为了确保正确性，控制逻辑必须确保每个输出总线最多连接到一个输入总线（以防止输入数据混合）。
然而，为了提高性能，逻辑还必须最大化并行通信的线路卡对数量。虽然当所有N个输出总线同时忙碌时
可以实现理想的并行性，但在实际情况中，并行性受到两个因素的限制。首先，某些输出线路卡可能没
有数据。其次，两个或更多输入线路卡可能希望向同一个输出线路卡发送数据。由于一次只有一个输入
可以胜出，这就限制了数据吞吐量，如果其他“失败”的输入无法发送数据的话。
因此，尽管有广泛的并行性，但主要的竞争发生在输出端口。如何在最大化并行性的同时解决输出端口
的竞争问题呢？DEC的GigaSwitch首次发明并使用了一种简单而优雅的调度方案来解决这个问题。图5
展示了其中使用的所谓“取票”算法的操作示例。
其基本思想是，每个输出线路卡S本质上为所有等待向S发送数据的输入线路卡R维护一个分布式队列。S
的队列实际上使用一种类似于熟食店柜台的简单票号机制存储在输入线路卡本身（而不是在S处）。如果
线路卡R想要向线路卡S发送一个数据包，它首先通过一条单独的控制总线向S发出请求；然后S通过控制
总线向R返回一个队列号。队列号是R在S的输出队列中的位置编号。
然后R监控控制总线；每当S完成接收一个新数据包时，S会在控制总线上发送它当前正在服务的队列
号。当R注意到它的号码正在被“服务”时，R将其数据包放置在R的输入数据总线上。同时，S确保R - S
交叉点导通。
为了更直观地理解这个算法，考虑图5，在顶部框架中，输入线路卡A有三个数据包，分别发往输出端口
1、2和3。B也有三个类似的数据包发往相同的输出端口，而C有数据包发往输出端口1、3和4。在这个例
子中，假设数据包大小相同（尽管这对于取票算法的运行不是必需的）。




                                                    4
图5 在取票调度机制中，所有输入端口都有一个单一的输入队列，队列中的每个数据包都标有其目的地
输出端口号。因此，在顶部框架中，输入A、B和C向输出端口1发送请求。输出端口1（顶部中间）将第一
个号码分配给A，第二个号码分配给B，依此类推，这些号码用于对输出端口的访问进行序列化
每个输入端口仅处理其队列头部的数据包。因此，算法开始时，每个输入端口通过控制总线向其输入队
列头部数据包的目的地输出端口发送请求。因此，在这个例子中，每个输入端口都向输出端口1发送请求
以获得发送数据包的许可。
一般来说，票号是一个小整数，会循环使用，并按到达顺序发放，就像在熟食店一样。在这种情况下，
假设A的请求首先到达串行控制总线，其次是B，然后是C，尽管左上角的图表看起来这些请求是同时发
送的。由于输出端口1一次只能处理一个数据包，它对请求进行序列化，将T1返回给A，将T2返回给B，
将T3返回给C。
因此，在顶部一行中间的图表中，输出端口1还在另一条控制总线上广播它当前正在服务的票号（T1）。
当A看到它有一个与输入1匹配的号码时，在顶部右边的图表中，A然后将其输入总线连接到1的输出总
线，并在其输入总线上发送其数据包。因此，在最上面一行图表结束时，A已经将其输入队列头部的数据
包发送到了输出端口1。不幸的是，所有其他输入端口，B和C，都被困在等待从输出端口1获得匹配的票
号。
图5中第二行开始时，A为其输入队列中现在位于头部的数据包向输出端口2发送请求；A收到端口2的票
号T1。在第二行中间的图表中，端口1宣布它准备好处理T2，端口2宣布它准备好处理票号T1。结果在第

                                                5
二行最右边的图表中，A连接到端口2，B连接到端口1，相应的数据包被传输。
第三行图表的开始类似，A和B分别向端口3和2发送请求。请注意，可怜的C仍然被困在等待它在两轮前
获得的票号T3，该票号将由输出端口1宣布。因此，在其队列头部的数据包得到服务之前，C不再发出更
多请求。A收到端口3的T1，B收到端口2的T2。然后端口1广播T3（终于！），端口2广播T2，端口3广播
T1。这导致了第三行最后一个图表的结果，其中CrossBar交换机连接A和3、B和2以及C和1。
取票方案在控制状态方面具有良好的扩展性，每个输出端口仅需两个 log2 N 位的计数器，用于存储当
前正在服务的票号以及已分配的最高票号。这使得DEC的千兆交换机（GigaSwitch）能够轻松扩展到36
                                    ​




个端口（Souza等人，1994），即便在20世纪90年代初期，片上内存有限的情况下也是如此。该方案采
用分布式调度器，每个输出端口的仲裁由每个线路卡上的一个所谓的GPI芯片完成；GPI芯片通过控制总
线进行通信。
GPI芯片必须对（串行）控制总线进行仲裁，以便向输出线路卡发出请求并获取队列号。由于控制总线是
广播总线，输入端口可以通过观察排在它前面的端口的服务情况，来判断自己的轮次，然后指示
CrossBar交换结构建立连接。
除了控制状态简单之外，取票方案的优点还在于能够直接处理可变大小的数据包。输出端口在完成当前
数据包的接收后，可以异步广播下一个票号；不同的输出端口可以在任意时间广播它们当前的票号。因
此，与后面描述的所有其他方案不同，该方案无需将数据包拆分成“信元”并在之后重新组装，从而避免
了报头和控制开销。另一方面，由于队首阻塞（HOL blocking）的存在，取票方案的并行性受到限制，
下一节我们将探讨这一现象。
取票方案还具备一个称为搜寻组（hunt groups）的优良特性。任意一组线路卡（并非仅仅是物理上相邻
的线路卡）都可以聚合在一起，形成一个带宽更高的有效链路，即搜寻组。例如，三条100Mbps的链路
可以聚合起来，等效于一条300Mbps的链路。
搜寻组的概念仅需对原始调度算法进行微小修改，因为组内的每个GPI芯片都可以在控制总线上观察彼此
的消息，从而使（公共）票号的本地副本保持一致。发往该组的下一个数据包由搜寻组中第一个空闲的
输出端口处理，这与有多个服务员的熟食店的运作方式类似。虽然基本的搜寻组可能会导致发往不同链
路的数据包重新排序，但通过一个简单的确定性哈希函数进行微小修改，可使来自一个输入的数据包仅
发送到搜寻组中的一个输出端口。这种修改避免了数据包重新排序，但代价是并行性降低。
由于千兆交换机是一种网桥设备，它必须处理局域网多播。由于取票调度机制通过每个输出端口的单独
GPI芯片进行分布式调度，很难协调所有调度器以确保每个输出端口都空闲。此外，等待所有端口为多播
数据包提供空闲票号，会导致一些本可提前处理该数据包的端口被阻塞，从而浪费吞吐量。因此，多播
由中央处理器通过软件处理，因此被视为“二等”业务。
队首阻塞
先不考虑图5的内部机制，注意在三次迭代（三个输入端口，三次迭代）中，原本有九次潜在的传输机
会。然而，在右下角图中所示的连接完成后，B的队列中还有一个数据包，C的队列中有两个数据包。因
                                                   6
此，原本可能传输的九个数据包中，实际上只传输了六个，这使得并行性的优势没有得到充分发挥。
图6简要描绘了仅关注输入 - 输出行为的情况。该图展示了每个输出端口在每个数据包传输时刻所发送
的数据包。每个输出端口都有一个相关的时间轴，上面标注着在相应时间段内发送数据包的输入端口，
如果没有则留空。需要注意的是，这张图延续了图5中的示例，又进行了三次迭代，直到所有输入队列都
为空。




图6 取票式等方案导致的队首阻塞示例。对于每个输出端口，绘制了一个水平时间轴，标注了在相应时
间段内将数据包发送到该输出端口的输入端口，若没有则标记为空白。注意大量的空白，这表明可能存
在浪费的机会，限制了并行性
从图6右侧的图表中很容易看出，只有大约一半的传输机会（更准确地说，24次中有9次）被利用了。当
然，对于某些场景，没有算法能做得更好。然而，其他算法，如iSLIP，可以挖掘更多的并行机会，并且
能在四次迭代内完成同样九个数据包的传输，而不是六次。
在图6的第一次迭代中，所有输入端口都有数据包等待发往输出端口1。由于一次只能有一个输入端口
（即A）向输出端口1发送数据包，B（和C）的整个队列都被阻塞，只能等待A完成传输。由于整个队列
的进度取决于队列头部数据包的传输情况，这种现象被称为队首阻塞（Head-of-line blocking，HOL
blocking）。iSLIP和PIM通过允许被阻塞数据包后面的数据包继续传输（例如，在图6的第一次迭代中，
B输入队列中发往输出端口2的数据包可以被发送到输出端口2）来克服这一限制，但代价是调度算法更
加复杂。
可以使用一个简单的均匀流量模型来分析队首阻塞导致的吞吐量损失。假设每个输入队列的队首有一个
数据包，该数据包发往N个输出端口中每个端口的概率为(1/N)。因此，如果两个或更多输入端口将数据
包发往同一个输出端口，除了一个输入端口外，其他所有输入端口都将被阻塞。由于队首阻塞，其他输
入端口的整体吞吐量就 “损失” 了。
更准确地说，假设数据包大小相等，并且在一次初始试验中，随机过程在每个输入端口从1到N中均匀选
择一个目的端口。我们不关注输入端口，而是关注输出端口o空闲的概率。这仅仅是N个输入端口中没有
一个选择o的概率。由于每个输入端口不选择o的概率为(1 - 1/N)，那么所有N个输入端口都不选择o的概
率为((1 - 1/N)^N) 。这个表达式迅速收敛到(1/e) 。因此，o忙碌的概率是(1 - 1/e)，约为0.63。因此，


                                                               7
交换机的吞吐量不是(N*B)（如果所有N个输出链路都以B比特/秒的速率忙碌运行，理论上可以达到这个
值），而是这个最大值的63%，因为37%的链路处于空闲状态。
这个分析比较简单，并且（错误地）假设每次迭代是相互独立的。实际上，在一次迭代中未发送的数据
包必须在下一次迭代中尝试发送（无需再次随机选择目的地）。经典分析（Karol等人，1987）去除了独
立试验的假设，结果表明实际利用率略低，更接近58.6%。
但是，均匀流量分布现实吗？显然，这种分析非常依赖于流量分布，因为如果所有流量都发往一个服务
器端口，那么任何交换机都无法表现良好。简单分析表明，通过使用搜寻组、在CrossBar交换结构中相
对于链路采用加速技术，以及假设更符合实际的流量分布（即多个客户端向少数服务器发送流量），可
以降低队首阻塞的影响。
然而，应该清楚的是，确实存在一些流量分布情况，队首阻塞会对吞吐量造成极大的损害。想象一下，
每个输入链路都有B个数据包发往端口1，接着是B个数据包发往端口2，依此类推，最后是B个数据包发
往端口N。所有输入端口的输入数据包分布都是如此。因此，显然在调度发往端口1的初始数据包组时，
由于队首阻塞，交换机实际上每次只能从每个输入端口发送一个数据包。因此，如果B足够大，交换机的
吞吐量将降至其可能吞吐量的(1/N) 。另一方面，我们将看到，使用虚拟输出队列（VOQ，本章后面会定
义）的交换机在相同情况下可以实现接近100%的吞吐量。这是因为在这种方案中，每组B个数据包在每
个输入端口都分别存储在不同的队列中。
通过输出排队避免队首阻塞
队首阻塞问题被发现后，有大量论文提出用输出排队取代输入排队来避免队首阻塞。假设数据包可以在
不经过输入排队的情况下被发送到输出端口，那么发往繁忙输出端口的数据包P就不可能阻塞其后面的数
据包。这是因为数据包P被发送到输出端口的队列中，在那里它只会阻塞发往同一输出端口的其他数据
包。
实现这一点的最简单方法是使交换结构的运行速度比输入链路快N倍。这样，即使在给定的时隙内所有N
个输入都将数据包发送到同一个输出端口，所有N个信元也可以通过交换结构发送到输出端进行排队。因
此，纯粹的输出排队要求交换结构内部有N倍的速度提升。这可能成本高昂或难以实现。
淘汰开关（knockout-switch）设计（Yeh等人，1987）提供了一种实际的输出排队实现方式。假设在任
何时隙内接收到N个发往同一目的地的信元的情况很少见，并且预期数量为k，k远小于N。那么，可以通
过将交换结构链路的运行速度设计为输入链路的k倍（而不是N倍）来优化预期情况（P11）。这在交换结
构内部实现了大幅节省。可以通过硬件并行性（P5）使用k条并行总线来实现这一点。
与取票式方案不同，本章中其余的所有方案，包括淘汰开关方案，都依赖于将数据包拆分成固定大小的
信元。在本章的其余部分，将使用信元代替数据包，需要始终记住，在数据包进入交换结构之前，必须
有一个初始阶段将数据包拆分成信元，并在输出端口重新组装。
除了更快的交换速度，淘汰开关方案还需要一个输出队列，该队列接收信元的速度比输出链路的速度快k
倍。使用简单规格设计的幼稚方案可以使用快速先进先出（FIFO）队列来实现，但成本高昂。而且，更
                                                       8
快的FIFO有些过度，因为显然缓冲区无法长期维持其输入和输出速度之间的不平衡。因此，可以放宽缓
冲区规格（P3），使其仅处理信元到达速度比取出速度快k倍的短时段。这可以通过内存交错和k个并行
内存来处理。一个分配器（其运行速度确实为k倍）以循环顺序将到达的信元分发到k个内存中，并且以
相同的顺序读出离开的信元。
最后，该设计必须公平地处理预期情况未满足的情况，即(N>k)个信元同时被发送到同一个输出端口。理
解通用解决方案的最简单方法是先理解三个更简单的情况。
 ● 两个竞争者，一个获胜者：在最简单的情况下，(k = 1)且(N = 2)，仲裁器必须从两个选项中公平地
   选择一个信元。这可以通过构建一个基本的2×2交换元件（称为集中器）来实现，该集中器随机选择
   一个获胜者和一个失败者。获胜者输出的就是被选中的信元。失败者输出在一般情况下很有用，在
   这种情况下，会组合使用多个基本的2×2集中器。
 ● 多个竞争者，一个获胜者：现在，考虑(k = 1)（一次只能接受一个信元）且(N>2)（仲裁器必须从多
   个信元中公平选择）的情况。一种简单的策略是使用分治法（P15，高效数据结构）创建一个由2×2
   集中器组成的淘汰树。就像在网球锦标赛的第一轮比赛中一样，使用(N/2)个基本的2×2集中器对信
   元进行配对，每个集中器代表一场比赛。这构成了树的底层。第一轮的获胜者被发送到第二轮的
   (N/4)个集中器，依此类推。“锦标赛” 以最后一轮结束，根集中器在这一轮中选择一个获胜者。请
   注意，每个集中器的失败者输出仍然被忽略。
 ● 多个竞争者，多个获胜者：最后，考虑一般情况，即对于任意的k和N值，必须从N个可能的信元中
   选择k个信元。一个简单的想法是创建k个单独的淘汰树来计算前k个获胜者。然而，为了公平起见，
   前面树中的失败者必须被发送到后续位置的淘汰树中。这就是为什么基本的2×2淘汰集中器有两个输
   出，一个用于获胜者，一个用于失败者，而不仅仅是一个用于获胜者。失败者输出被路由到后续位
   置的树中。
请注意，如果要从八个选项中选择四个信元，最简单的设计是将八个选项（两两分组）分配给四个2×2淘
汰集中器。这种逻辑会选出四个获胜者，即所需的数量。虽然这种逻辑显然比使用四个单独的淘汰树简
单得多，但可能非常不公平。例如，假设两个流量很大的源(S1)和(S2)恰好被分在一组，而另一个流量较
大的源(S3)与一个流量较小的源分在一组。在这种情况下，(S1)和(S2)获得的流量大约是(S3)的一半。为
了避免这些不公平的情况，淘汰逻辑使用k个单独的树，每个位置对应一个树。
实现这些树的简单方法是在完成位置(j - 1)树的所有逻辑后，严格开始运行位置j树的逻辑。这确保了收
集到所有符合条件的失败者。练习中探讨了一种更快的实现技巧。希望这个设计能让你相信，实现公平
是很困难的。这也是在讨论iSLIP时将再次探讨的主题。
虽然淘汰开关因其引入的技术而值得深入理解，但它实现起来很复杂，并且对流量分布有一定假设。对
于实际拓扑结构而言，这些假设并不成立，在实际情况中，经常会有超过k个客户端同时向一个热门服务
器发送数据。更重要的是，研究人员设计出了相对简单的方法来应对队首阻塞，而无需采用输出排队。
通过虚拟输出排队避免队首阻塞

                                                     9
一种经受住时间考验的解决方案是虚拟输出排队（Virtual Output Queuing，VOQ）（Tamir和Frazier，
1988）。事实上，它比上一节中描述的大多数输出排队解决方案提出的时间更早。VOQ方法是重新考虑
输入排队，并对其进行改进以避免队首阻塞。它通过允许输入端口不仅调度其输入队列的头部数据包，
还调度其他数据包来实现这一点，当队列头部数据包被阻塞时，其他数据包可以继续传输。乍一看，这
似乎非常困难。每个队列中可能有数十万个信元；为每个信元维护哪怕1位的调度状态都需要大量的内存
来存储和处理。
然而，第一个重要的发现是，每个输入端口队列中的信元只能发往N个可能的输出端口。假设在输入队列
X中，信元(P1)在信元(P2)之前，并且(P1)和(P2)都发往同一个输出队列Y。那么，为了保持先进先出
（FIFO）的行为，无论如何(P1)都必须在(P2)之前被调度。因此，在(P1)完成调度之前尝试调度(P2)显然
是浪费资源。因此，可以通过不调度发往每个不同输出端口的第一个信元之外的其他信元来避免明显的
浪费（P1）。
VOQ方法的核心思想是利用一种自由度（P13），将图5中的单个输入队列分解为图7中每个输入端口对
应每个输出端口的单独输入队列。这些队列被称为虚拟输出队列。请注意，图7左上角的图包含与图6相
同的输入信元，只是现在它们被放置在单独的VOQ中。
随着VOQ概念的引入，我们在交换过程中面临一个新的计算问题。对于具有单个（组合）数据包队列的
输入端口，它应该始终（如果可能的话）与队列头部数据包的目的输出端口配对；而对于具有N个VOQ
的输入端口，它有多达N个配对选项，并且必须从中选择一个。更一般地说，具有N个输入端口和N个输
出端口的交换机必须从其(N^2)个VOQ（每个输入端口有N个）中决定在一个时隙内为哪些（或更少）
VOQ提供服务。
正如我们接下来将解释的，在CrossBar交换机中，做出这样的配对决策涉及计算二分匹配。自20世纪70
年代以来，人们就知道，一般来说，二分匹配是一个非平凡的计算问题。实际上，对于大型快速交换机
（即N很大且时隙很短的情况），通用的二分匹配算法计算成本太高。这可能就是为什么VOQ方法在首
次提出后大约十年才被广泛采用，直到专门为交换设计的计算效率高的二分匹配算法，如PIM
（Anderson等人，1993）和iSLIP（McKeown，1999）被发现。

输入排队交换作为二分匹配问题
我们可以将CrossBar交换机的调度问题表述为一个二分匹配问题，具体如下。一个(N×N)的输入排队
CrossBar交换机可以建模为一个加权完全二分图，其中两个不相交的顶点集分别是N个输入端口和N个输
出端口。在这个二分图中，任意输入端口i和任意输出端口j之间都有一条边，这条边对应于输入端口i处用
于缓冲发往输出端口j的数据包的VOQ。这条边的权重定义为该VOQ的长度（即缓冲的信元数量）。一组
这样的边构成一个有效的CrossBar调度，即一个匹配，如果其中任意两条边不共享一个公共顶点。一个
匹配的权重是属于它的所有边的总权重（即所有相应VOQ的总长度）。


                                                              10
在CrossBar调度问题中，有三种类型的匹配起着重要作用：（1）极大匹配；（2）最大匹配；（3）最
大权重匹配（MWM）。如果一个匹配S在添加任何权重不为零且不在S中的边后不再是匹配，那么这个
匹配S就被称为极大匹配。具有最多非零权重边数的匹配称为最大匹配或最大基数匹配。极大匹配和最大
匹配都不考虑边的权重，而最大权重匹配会考虑。最大权重匹配是所有匹配中总权重最大的匹配。在交
换的背景下（其中所有边的权重都是非负的），任何最大匹配或最大权重匹配也是极大匹配，但反之通
常不成立。
虽然使用最大权重匹配（在每个时隙）进行CrossBar调度可以证明在所有流量模式下都能保证100%的吞
吐量（McKeown等人，1999；Tassiulas和Ephremides，1992），并且在经验上可以实现非常低的排队
延迟（McKeown等人，1999），但目前最先进的串行最大权重匹配算法具有非常高的计算复杂度，为
(O(N{2.5} log W))（Duan和Su，2012），其中W是边（VOQ）的最大可能权重（长度）。同样地，最大匹配的情况更不理想：它的时间
复杂度略低，为(O(N{2.5}))（Hopcroft和Karp，1973），然而使用最大匹配作为CrossBar调度通常不能在所
有流量模式下保证100%的吞吐量。长期以来，极大匹配族被认为是CrossBar调度中一种特殊的具有成本
效益的匹配族。一方面，存在高效的分布式算法来计算极大匹配。我们将在下两节中描述两种这样的算
法，即并行迭代匹配（PIM）（Anderson等人，1993）和iSLIP。另一方面，使用极大匹配作为
CrossBar调度可以证明在所有流量模式下至少能实现50%的吞吐量（Dai和Prabhakar，2000），并且
像PIM和iSLIP这样的极大匹配算法通常在实际应用中能提供更好的吞吐量性能。总之，众所周知，除了
在某些良性流量模式（如均匀流量）下，PIM和iSLIP都无法实现100%的吞吐量。
并行迭代匹配（PIM）
严格来说，PIM是一种近似极大匹配算法，也就是说它只是以高概率输出一个极大匹配，iSLIP也是如
此。PIM是一种随机算法（P3a）。具体而言，它可以看作是对Israel和Itai（1986）提出的随机（近似）
极大匹配算法在交换领域的改编。在PIM中，任何输入端口的调度需求只需一个N位的位图，其中N是交
换机的规模。在位图中，第i位为1表示至少有一个信元发往输出端口i。因此，如果N个输入端口中的每个
端口都通过一个N位向量来描述其调度需求，那么调度器只需处理(N^{2})位信息。对于较小的(N ≤32)
，通过控制总线传输或存储在控制内存中的这些位数并不多。




                                                                        11
图7 并行迭代匹配（PIM）方案的工作原理是，让所有输入端口并行地向其希望到达的所有输出端口发送
请求。然后，PIM使用随机化方法进行公平匹配，使得收到多个请求的输出端口随机选择一个输入端口，
而收到多个授权的输入端口则随机选择一个输出端口进行发送。这个三轮过程可以迭代进行，以扩大匹
配的规模。
在图7左上角的图中，从每个输入端口向其非空虚拟输出队列（VOQ）对应的输出端口发送一条线，以此
表示请求的传递。注意，A没有向端口4发送线，因为它没有发往端口4的信元。还要注意，输入端口C为
其输入端口C中发往输出端口4的信元发送了请求，而在图5的单输入队列场景中，该信元是输入队列C中
的最后一个信元。
目前仍需要一种调度算法，将资源与需求进行匹配。虽然调度算法很巧妙，但作者认为真正的突破在于
认识到，使用VOQ后，输入队列调度可以避免队首阻塞（HOL blocking）。为了使图7中的示例与图5相
对应，假设图5场景中的每个数据包在图7中都被转换为单个信元。
为了说明PIM中使用的调度算法，观察图7左上角的图，输出端口1收到了来自A、B和C的三个请求，但在
下一个时隙只能处理一个。一种简单的请求选择方法是随机选择（P3a）。因此，在授权（Grant）阶段
（图7中间上方的图），输出端口1随机选择了B。同样，假设端口2从A和B中随机选择A，端口3从A、B
和C中随机选择A，最后端口4选择其唯一的请求者C。
然而，解决输出端口的竞争是不够的，因为还存在输入端口的竞争。两个输出端口可能会随机授权给同
一个输入端口，而该输入端口必须准确选择一个输出端口来发送信元。例如，在图7中间上方的图中，A

                                                   12
收到了来自输出端口2和3的授权，面临选择难题。因此，需要第三个阶段——接受（Accept）阶段，在
这个阶段，每个输入端口随机选择一个输出端口（既然在授权阶段使用了随机化，为什么不再次使用
呢？）。
因此，在图7右上角的图中，A随机选择了端口2。B和C别无选择，分别选择了端口1和4。CrossBar交换
结构建立连接，数据包从A发送到端口2，从B发送到端口1，从C发送到端口4。在这种情况下，找到的相
应匹配恰好是一个极大匹配，但在某些情况下，随机选择可能会导致匹配不是极大匹配。例如，在不太
可能出现的情况下，端口1、2和3都选择A，那么匹配的规模就只有2。
在这种情况下（图中未显示），算法可能值得屏蔽掉所有已匹配的输入和输出，并对同一即将到来的时
隙进行更多次迭代。如果当前迭代的匹配不是极大匹配，进一步的迭代将使匹配规模至少增加1。请注
意，后续迭代不会使匹配变差，因为现有匹配在迭代过程中会被保留。虽然对于N个输入，达到极大匹配
的最坏情况时间是N次迭代，但简单的推理表明，预期的匹配次数更接近log N。DEC AN - 2的实现
（Anderson等人，1993）在一个30端口的交换机上使用了三次迭代。
然而，我们在图7中的示例每次匹配只使用一次迭代。中间一行显示了第二个时隙的第二次匹配，例如，
A和C都请求端口1（B没有请求，因为B发往1的信元在上一个时隙已发送）。端口1随机选择C，最终的匹
配结果是A与3、B与2、C与1。第三行显示了第三次匹配，这次匹配规模为2。在第三次匹配结束时，只
有输入队列B中发往端口3的信元未发送。因此，在四个时隙内（其中第四个时隙使用较少，本可用于发
送更多流量），所有流量都被发送出去。这显然比图5中的取票式示例更高效。
用iSLIP避免随机化
并行迭代匹配是一种具有开创性的方案，因为它提出了一个观点：通过巧妙的并行化，可以以合理的计
算和硬件成本计算出相当不错的二分匹配。正如罗杰·班尼斯特（Roger Bannister）首次在四分钟内跑
完一英里后，其他人可以在此基础上进一步改进一样。但是，PIM存在两个潜在问题。第一，它使用随机
化，在非常高的速度下可能很难生成合理的随机数源。第二，它需要对数次迭代才能达到极大匹配。鉴
于对数次迭代中的每一次都需要三个阶段，并且整个匹配决策必须在最短的数据包到达时间内完成，因
此，如果有一种匹配方案能够在一两次迭代内接近极大匹配，那就更好了。
iSLIP是一种非常流行且有影响力的方案，它本质上是对PIM进行“去随机化”，并且在一两次迭代后就能
实现非常接近极大匹配的结果。其基本思想非常简单。在PIM中，当一个输入端口或输出端口收到多个请
求时，为了公平起见，它会从这些请求中随机选择一个“获胜”请求。以太网通过随机化来提供公平性，
而令牌环则使用由旋转令牌实现的轮询指针来实现公平性。
同样，iSLIP通过使用旋转指针以轮询方式在多个竞争者中选择下一个获胜者来提供公平性。虽然轮询指
针在初始时可能会同步并导致类似队首阻塞的情况，但从长期来看，至少在模拟中，它们会逐渐摆脱同
步并实现极大匹配。因此，iSLIP的精妙之处不在于使用轮询指针，而在于N个这样的指针在并发运行时
明显没有长期同步。


                                                   13
更准确地说，每个输出（分别对应输入）端口维护一个指针g（分别对应a），初始时设置为第一个输入
（分别对应输出）端口。当一个输出端口需要在多个输入请求中进行选择时，它会选择大于或等于g的最
小输入端口号。同样，当一个输入端口需要在多个输出端口请求中进行选择时，它会选择大于或等于a的
最小输出端口号，其中a是该输入端口的指针。如果一个输出端口与输入端口X匹配，那么该输出端口的
指针会按循环顺序递增到大于X的第一个端口号（即g变为(X + 1)，除非X是最后一个端口，在这种情况下
g会回绕到第一个端口号）。
这种“旋转优先级”的简单机制确保了每个资源（输出端口、输入端口）在所有竞争者之间得到合理公平
的共享，代价是除了每次迭代所需的(N^{2})个调度状态外，还需要额外的(2N)个(\log _{2} N)指针。
图8和图9展示了与图7（以及图5）相同的场景，但使用了两次迭代的iSLIP。由于每一行代表一次匹配的
一次迭代，所以每次匹配用两行表示。因此，图8的三行展示了该场景的前1.5次匹配。同样，图9展示了
剩余的1.5次匹配。
图8左上角的图与图7相同，每个输入端口向其有信元发往的输出端口发送请求。然而，一个不同之处在
于，每个输出端口都有一个所谓的授权指针g，所有输出端口的g初始值都设为A。同样，每个输入端口都
有一个所谓的接受指针a，所有输入端口的a初始值都设为1。




图8 一个iSLIP示例场景的1.5轮匹配
iSLIP的确定性在授权阶段立即产生了差异。比较图8中间上方和图7中间上方的图。例如，当输出端口1
收到来自所有三个输入端口的请求时，它授予A，因为A是大于或等于(g_{1}=A)的最小输入端口。相比之
                                                       14
下，在图7中，端口1随机选择输入端口B。在这个阶段，iSLIP的确定性似乎是一个真正的劣势，因为A向
输出端口3和4都发送了请求。由于3和4的授权指针(g_{3}=g_{4}=A)，端口3和4也都授予A，忽略了B和
C的请求。和之前一样，由于C是端口4的唯一请求者，C获得了来自4的授权。
当备受青睐的A从端口1、2和3收到三个授权时，A接受端口1。这是因为端口1是大于或等于(A)的接受指
针(a_{A})（其值为1）的第一个输出端口。同样，C选择4。完成选择后，A将(a_{A})递增到2，C将
(a_{C})递增到1（在循环顺序中，1大于4的下一个数是1）。只有在这个阶段，输出端口1才将其授权指针
(g_{1})递增到B（比上一次成功授权的端口号大1），端口4同样将其递增到A（比C大1，在循环顺序
中）。
请注意，虽然端口2和3向A授予了授权，但它们不会递增其授权指针，因为A拒绝了它们的授权。如果它
们递增了指针，就可能会出现一种情况：输出端口在未成功授权后不断将其授权指针递增到某个输入端
口I之后，从而持续使输入端口I处于饥饿状态。还要注意，此时的匹配规模只有2；因此，与图7不同，这
个iSLIP场景可以通过第二次迭代得到改进，如图8的第二行所示。注意，在第一次迭代结束时，匹配的
输入和输出并没有像图7右上角那样用实线连接（表示数据传输）。数据传输将等待最后一次（在这种情
况下是第二次）迭代结束。
第二次迭代（图8的中间一行）仅从不匹配的输入端口（即B）向之前未匹配的输出端口发送请求。因
此，B向2和3发送请求（尽管B也有一个发往1的信元，但这次不向1发送请求）。2和3都授予B，B选择2
（这是大于或等于其接受指针1的最小端口）。有人可能认为B应该将其接受指针递增到3（1加上最后接
受的端口号2）。然而，为了避免饥饿，iSLIP在除第一次迭代外的其他迭代中不递增指针，原因将在后
面解释。
因此，即使B连接到2，2的授权指针仍然保持在A，B的接受指针仍然保持在1。由于这是最后一次迭代，
所有匹配对，包括之前迭代中匹配的对（如A与1），都进行连接并进行数据传输（实线表示）。
第三行展示了授权和接受指针的初始同步是如何被打破的。因为只有一个输出端口授予了A，该端口（即
1）得以继续进行，并在这次为A之后的端口提供优先级。因此，即使A有第二个发往1的数据包（在这个
例子中没有），1仍然会授予B。
读者应该仔细研究图8和图9中的其余行，检查授权和接受指针的更新规则，并查看每一轮中哪些数据包
被交换。底线是，到图9第三行结束时，只剩下B发往3的信元需要交换。这显然可以在第四个时隙完
成。




                                                      15
图9 图8所示iSLIP示例场景的最后1.5轮匹配




图10 iSLIP在图6场景中如何避免队首阻塞以提高吞吐量
还要注意，当我们比较图8和图7时，iSLIP看起来比PIM差，因为iSLIP每次匹配需要两次迭代才能达到
PIM一次迭代所达到的匹配规模。然而，这更多地说明了iSLIP的启动代价，而非长期代价。在实际应用
中，一旦iSLIP指针不同步，iSLIP只需一次迭代就能表现得很好，一些商业实现也只使用一次迭代：
iSLIP非常受欢迎。
有人可能会将iSLIP总结为用输入和输出指针的轮询调度取代随机化的PIM。然而，这种描述忽略了iSLIP
的两个微妙之处。
  ● 授权指针仅在第三次阶段（即授权被接受后）递增：直观地说，如果输出端口o授予输入端口I，不能


                                                  16
   保证I会接受。因此，如果o在I之前递增其授权指针，可能会导致从I到o的流量持续饥饿。更糟糕的
   是，McKeown等人（1997）表明，这种简单的轮询方案会将吞吐量降低到仅63%（对于伯努利到
   达的流量），因为指针倾向于同步并同步移动。
 ● 所有指针仅在第一次迭代接受被授予后递增：同样，这个规则是为了防止饥饿，但具体情况更为微
   妙，练习中会要求你思考这个问题。
因此，iSLIP中除第一次迭代外的其他迭代中的匹配被视为一种“额外奖励”，可以提高吞吐量，而不会计
入端口的配额。
iSLIP实现说明
iSLIP硬件实现的核心是一个仲裁器，它在N个请求（编码为位图）中进行选择，以找到第一个大于或等
于固定指针的请求。这就是第2章中提到的可编程优先级编码器；该章还介绍了一种高效实现方式，其速
度几乎与优先级编码器一样快。交换机调度可以通过为每个输出端口配置一个这样的授权仲裁器（用于
仲裁请求），以及为每个输入端口配置一个接受仲裁器（用于仲裁授权）来完成。优先级和多播功能是
通过在输入到达仲裁器之前添加一个过滤器，集成到基本结构中的；例如，优先级过滤器会将除最高优
先级请求之外的所有请求清零。
原则上，单播调度器可以为每个端口设计一个单独的芯片，但由于状态量足够小，可以由一个单独的调
度器芯片处理，该芯片通过控制线路与每个端口连接。此外，多播算法要求每个优先级级别都有一个共
享的多播指针，这也意味着需要一个集中式调度器。然而，集中化意味着从端口线路卡向中央调度器发
送请求和决策会存在延迟。
为了容忍这种延迟，调度器（Gupta和McKeown，1999b）对每个虚拟输出队列（VOQ）中的m个信元
（Tiny Tera中为8个）和每个多播队列中的n个信元（Tiny Tera中为5个）进行流水线处理。这意味着在
Tiny Tera中，每个线路卡必须为每个单播VOQ传输3位信息，表示VOQ的大小，最大为8。对于32个输
出端口和四个优先级级别，每个输入端口必须发送384位的单播信息。每个线路卡还需要为每个优先级级
别中的五个多播数据包的每个扇出（每个扇出32位）进行通信，这导致总共需要传输640位。32×1024
位的输入信息存储在片上静态随机存取存储器（SRAM）中。然而，为了实现更高的速度，每个队列头部
的信息（较小的状态，例如每个单播VOQ仅1位）存储在速度更快但密度较低的触发器中。
现在考虑处理多次迭代的情况。注意，请求阶段仅在第一次迭代时出现，并且每次迭代只需屏蔽已匹配
的输入即可进行修改。因此，K次迭代似乎至少需要2K个时间步长，因为每次迭代的授权和接受步骤各
需要一个时间步长。乍一看，架构似乎规定在第k次迭代的接受阶段之后，才开始第k + 1次迭代的授权阶
段。这是因为需要知道输入端口I在第k次迭代中是否被接受，以避免在第k + 1次迭代中对该输入进行授
权。
实现部分流水线处理的关键在于一个简单的观察（Gupta和McKeown，1999b）：如果输入端口I在第k
次迭代中收到任何授权，那么它必须恰好接受一个，因此在第k + 1次迭代中不可用。因此，可以放宽实
现规范（P3），允许第k + 1次迭代的授权阶段在第k次迭代的授权阶段之后立即开始，从而与第k次迭代

                                                      17
的接受阶段重叠。为此，我们只需使用第k次迭代结束时对输入端口I的所有授权的“或”运算结果，来屏蔽
第k + 1次迭代中I的所有请求。
这使得K次迭代的整体完成时间几乎减少了一半，从2k个时间步长减少到k + 1个。例如，Tiny Tera的
iSLIP实现（Gupta和McKeown，1999b）使用175MHz的时钟速度，在51纳秒（大约OC - 192速度）内
完成iSLIP的三次迭代；鉴于每个时钟周期大约为5.7纳秒，iSLIP大约有9个时钟周期用于完成操作。由于
每个授权和接受步骤需要两个时钟周期，因此流水线处理对于在9个时钟周期内完成三次迭代至关重要；
如果采用简单的迭代技术，则至少需要12个时钟周期。
通过学习计算近似最优匹配
像最大权重匹配（MWM）这样的重匹配，通常在吞吐量和延迟性能方面都是很好的匹配。在大多数用于
交换的二分匹配算法中，如前面描述的PIM和iSLIP，每个时隙的CrossBar调度（匹配）计算都是从头开
始的，换句话说，忽略了之前的计算结果。这在计算上可能显得很浪费（P1）。例如，前一时隙计算出
的重匹配在当前时隙通常仍然是重匹配，因为在前一时隙中，最多可能有N个数据包离开，但在重载情况
下，可能会有接近N个数据包到达。直观地说，如果我们能够从之前时隙使用的匹配中学习到重边，那么
我们可能能够以比从头计算低得多的时间复杂度，增量式地（P12a）计算出新的重匹配。实际上，我们
可以通过接下来几节将描述的一系列自适应算法来实现这一点。
除了前一时隙使用的二分匹配之外，自适应算法还可以从早期计算中学习其他内容，以帮助更快地计算
出良好的新匹配。例如，一个智能的自适应算法可以在过去的许多时隙中逐渐发展并保持对所有(N^{2})
条边（VOQ）权重的 “态势感知”，这可以帮助算法做出快速而明智的匹配决策。在接下来的几节中，我
们假设输入端口和输出端口的数量均为N，并且允许边的权重为0（即相应的VOQ为空）。在这两个假设
下，最大匹配总是包含N条边，并且必然是完全匹配。因此，在本节中我们始终使用 “完全匹配” 这个术
语，以避免任何混淆。
抽样比较：一种极其简单的自适应算法
我们从Tassiulas（1998）提出的自适应算法开始介绍，我们将其称为抽样比较算法。这个算法非常简
单：在每个时隙t，抽样一个（均匀）随机匹配（P3a），记为(R(t))，将其权重与前一时隙使用的匹配
(S(t - 1))的权重进行比较，并将权重较大的匹配作为(S(t))（即用于当前时隙t）。与最大权重匹配算法相
比，抽样比较算法的计算复杂度要低得多，为(O(N))。令人惊讶的是，与最大权重匹配算法一样，抽样比
较算法可以证明在所有流量模式下都能实现100%的吞吐量。然而，在重载情况下，其延迟性能较差，这
可以通过它在重载下经历的排队动态来解释。
可以证明，当前（即在时间t）交换算法的归一化吞吐量大致等于(S(t))的权重与时间t时最大权重匹配
（MWM）权重的比值。特别是，对于任何像抽样比较算法这样要实现100%吞吐量的交换算法，(S(t))的
权重最终必须非常接近最大权重匹配的权重。然而，由于（均匀）随机匹配(R(t))只有极小的概率具有足
够大的权重来超过(S(t - 1))（(S(t - 1))已经相当大），在抽样比较算法下，(R(t)) “获得足够权重” 以接
近最大权重匹配的速度极其缓慢。因此，在很长一段时间内，(S(t))的权重可能比最大权重匹配的权重要
                                                             18
小得多，在此期间交换机的吞吐量远低于100%，结果是在重载情况下，边的权重（VOQ长度）变得非常
大。这就解释了抽样比较算法在延迟性能方面较差的原因。然而，当（几乎）所有边的权重都变得非常
大时，随机匹配(R(t))的权重开始有相当大的概率超过(S(t - 1))的权重，并且(S(t))的权重开始更快地接近
最大权重匹配的权重。这就解释了为什么抽样比较算法最终能够实现100%的吞吐量。
QPS实现
我们现在描述一种数据结构和算法，它能让输入端口以队列比例方式对VOQ进行抽样，并且在需要时
（为接受交换服务）移除任何VOQ的队首（HOL）数据包，这两个操作的（每个端口）计算复杂度均为
(O(1)) 。
该算法包含两个步骤。第一步，我们从当前输入端口队列中的所有数据包中随机均匀抽样一个数据包。
具体来说，如果输入端口所有(N)个VOQ中共有(m)个数据包，那么每个数据包被抽样的概率为(1/m) 。通
过这种均匀抽样，长度为(m_j)的第(j)个VOQ中，有一个数据包被抽样的概率为(m_j/m) ，这正是QPS的
行为。
假设这样抽样出了一个数据包。第二步的一部分工作是找出这个数据包属于哪个VOQ，以便输入端口能
向相应的输出端口提议匹配，并附上队列长度。然而，还需要做更多工作。由于所有交换算法都严格按
照先进先出（FIFO）顺序处理VOQ中的数据包，如果这个提议成功（即被输出端口接受），并且输入和
输出端口对最终成为最终匹配的一部分，那么就需要定位并处理这个VOQ的HOL数据包，该数据包可能
是也可能不是抽样出的那个数据包。因此，第二步的另一部分工作是定位这个VOQ的HOL数据包。
在详细介绍之前，我们列出这种数据结构还需要支持的另外两个基本操作。第一个操作是，任何新到达
的数据包都必须记录在数据结构中，以便在逻辑上 “添加到它所属的VOQ的末尾”。第二个操作是，当调
度算法最终决定将输入端口与一个不同于提议的输出端口配对时（这可能是因为提议被拒绝，或者最初
接受的提议被调度算法覆盖，例如在QPS - SERENA的情况下，在SERENA的合并操作期间），需要定
位并移除（新）相应VOQ的HOL数据包，以便接受交换服务。接下来将展示，这两个操作的复杂度均为
(O(1)) 。
QPS数据结构
我们展示了QPS提议策略的两个步骤，在任何输入端口都可以通过一个主数据结构和一个辅助数据结构
在(O(1))时间内完成，这两个数据结构对所有输入端口都是相同的。图13（A）和（B）展示了在单个输入
端口上，其第(j)个VOQ的HOL数据包被选择用于（交换）服务之前和之后的数据结构。图的上半部分和
下半部分分别展示了主数据结构和辅助数据结构。




                                                         19
图13 单个输入端口上QPS数据结构的操作示意图（改编自Gong等人，2017）
主数据结构是一个包含(N)个记录的数组，对应于输入端口的(N)个VOQ。每个记录(j)（即数组项(j)）都关
联一个链表，该链表按数据包到达顺序指向VOQ中排队的数据包，从HOL数据包开始。链表中的每个节
点包含两个编码为 “⟨字母⟩”（例如A）的指针；一个指针指向数据包缓冲区中实际的数据包（图中未显
示），另一个指针指向辅助数据结构中的相应项（例如项A），我们将其称为反向指针。
为简单起见，图13仅展示了对应VOQ (j)的记录(j) 。每个记录都包含一个头指针和一个尾指针，分别指向
链表的头节点和尾节点。头指针用于在(O(1))时间内定位和移除头节点（即HOL数据包）；它还用于定位
和替换辅助数据结构中对应HOL数据包的数组项。尾指针用于在(O(1))时间内将新到达的数据包插入
“VOQ的末尾”（即第一个基本操作）。
辅助数据结构。图13的下半部分展示了用于抽样的辅助数据结构。假设输入端口所有(N)个VOQ中总共排
队(m)个数据包。辅助数据结构只是一个包含(m)个项的数组，每个项都是一个指针，指向主数据结构中
(N)个链表之一中的一个不同（数据包）节点（例如节点A）。
随着时间推移，尽管数据包有到达和离开，但辅助数据结构始终占据数组中连续的一块区域，其边界由
一个头指针和一个尾指针标识，如图13下半部分所示。这种连续性使得可以在(O(1))时间内从数组中随机
均匀抽样一个项（数据包），这是QPS的关键步骤之一。因此，在数据包到达和离开时，需要保持这种
连续性。数据包到达的情况比较简单：对应新数据包的项插入当前尾指针位置之后，然后更新尾指针。
数据包离开的情况稍微复杂一些：如果离开的数据包在块中留下 “空洞”，则将尾项移动到该空洞位置，
并更新尾指针。
在数据包离开的情况下，主数据结构中被先前尾项（现在移动到 “填充空洞” 位置）指向的（数据包）节
点，其反向指针需要更新为先前空洞的偏移量，即先前尾项现在所在的位置。这显然是一个(O(1))的过
程。类似的过程可用于在(O(1))时间内支持第二个基本操作。
示例说明。为了了解主数据结构和辅助数据结构如何协同工作以实现QPS，考虑图13中的示例，其中数
据包A从辅助数据结构的(m)个数据包中被抽样出来。然而，它不是HOL数据包，因此检查其目的（输
出）端口（即VOQ标识符），结果发现是(j) 。通过访问主数据结构中对应VOQ (j)的第(j)个记录，发现
                                                    20
HOL数据包是数据包B。现在，输入端口提议与输出端口(j)匹配。如果提议被接受，并且输入端口最终与
输出端口(j)匹配，那么数据包B将在当前时隙发往（输出端口(j) ）。主数据结构中第(j)个记录的头指针将
更新为（指向）E，即新的HOL数据包。这些变化在图13（B）中有所体现。这些操作，包括查找HOL数
据包以及更新两个数据结构，都只需要(O(1))时间。
1.3 联合输入输出排队
20世纪90年代末到21世纪初，一种名为联合输入输出排队（Combined Input and Output Queueing，
CIOQ）的交换架构得到了研究（Prabhakar和McKeown，1997；Stoica和Zhang，1998；Chuang等
人，1999；Iyer等人，2002；Yang和Zheng，2003；Giaccone等人，2004；Firoozshahian等人，
2007）。CIOQ的提出是为了通过在输出端口执行数据包调度（下一章的主题）来提供服务质量
（Quality of Service，QoS）保证。在CIOQ中，数据包在输入端口和输出端口都进行排队，这也正是该
架构名称的由来。CIOQ的设计目标是，就输出端口的数据包调度而言，完全模拟输出排队的效果，同时
又无需承担输出排队所带来的高昂成本（输出端口需要N倍的速度提升）。
现在我们详细阐述 “完全模拟” 的概念。对于一个输出排队的交换机，为了提供某种QoS保证，每个输出
端口j需要在输出端口j处的N个虚拟输入队列（Virtual Input Queues，VIQs）之间执行某种数据包调度策
略P（所有输出端口通用），其中第i个VIQ（在输出端口j处）保存着所有从输入端口i到达的数据包。在
输出排队的交换机中，在任何数据包到达实例下，P完全决定了这些数据包的服务顺序，因为每个数据包
在到达时可以直接 “飞” 到其输出端口，不会有任何延迟，这是输出排队交换机所保证的。如果一个交换
机在任何数据包到达实例下，都能产生与输出排队交换机完全相同的数据包服务顺序，而不管两个交换
机在输出端口用于数据包调度时使用何种P，那么就称这个交换机完全模拟了输出排队。
CIOQ交换机与输入排队交换机有很大不同，前者计算的匹配与后者有很大差异。在CIOQ交换机中，在
任何给定时间（周期），输入端口和输出端口之间的匹配是根据P分配给(N^{2})个VOQ或VIQ的队首
（HOL）数据包的截止期限或优先级分数来决定的；这样计算出的匹配π必须是一个稳定匹配（Gale和
Shapley，1962），即不存在一对输入端口和输出端口，它们都（根据这些截止期限或优先级分数）更倾
向于相互配对，而不是当前在π下的配对。相比之下，前面几节中描述的用于输入排队交换机的
CrossBar调度算法与这样的P毫无关系，也不计算稳定匹配。
CIOQ技术从未得到广泛应用，可能是因为它在两个方面成本较高。第一，CIOQ调度算法通常具有较高
的时间复杂度，为(O(N{2}))（每周期），因为为了计算出稳定匹配，需要对这(N{2})个截止期限或优先级分数进行某种
排序（Gale和Shapley，1962）。第二，它们只能保证50%的吞吐量（Chuang等人，1999；Dai和
Prabhakar，2000）。因此，这里的成本在于，为了提供QoS保证（使用P），输出端口需要2倍的速度
提升。Firoozshahian等人（2007）和Prabhakar（2009）解释了CIOQ未得到广泛应用的一个有趣观
点。
CIOQ很可能永远不会被广泛采用，原因如下。如前所述，我们已经有了优秀的输入排队交换算法，如
SW-QPS，它可以以非常低的成本提供出色的吞吐量和延迟性能。我们可以简单地使用这样的算法进行

                                                                 21
数据包从输入端口到输出端口的交换，然后在输出端口执行P。当负载不是太高时，任何数据包在相应输
入端口的排队延迟（由于交换导致）应该很小，因此极有可能在P规定的截止期限前到达其目的输出端
口。这样，我们几乎可以获得CIOQ的全部好处（P提供的QoS保证），而无需承担其高昂的成本。
1.4 扩展到更大和更快的交换机
由于现有的交换方案在扩展到大量端口和更高端口速度时存在各种障碍。这两种方法都通过降低两种主
要交换成本中的一种或两种来实现可扩展性目标，这两种成本分别是交换电路的整体尺寸和二分匹配计
算的时间复杂度。第一种方法基于分治原理，通过使用比单一CrossBar交换结构更节省空间的交换结
构，如Clos交换结构）和Benes交换结构，来减小交换电路的整体尺寸。另一种方法称为负载均衡交换
（Load-Balanced Switching，LBS），它将二分匹配计算的算法成本几乎降为零。然而，它通过使用
两个大型CrossBar交换结构而非一个来实现这一点。
到目前为止，本章主要集中讨论了适用于构建32端口路由器的较小规模交换机。在撰写本文时，大多数
已部署的互联网路由器都属于这一类别，有时这是有充分理由的。例如，建筑布线规范往往会限制从一
个配线间可以服务的办公室数量。因此，位于配线间的局域网交换机采用较小的端口尺寸通常就可以满
足需求。
最近有三个趋势有利于大型交换机的设计：
 1. 密集波分复用（DWDM）：在核心网络中使用密集波分复用技术，有效地在光链路上捆绑多个波
    长，这将显著增加核心路由器必须交换的逻辑链路数量。
 2. 光纤到户：在不久的将来，家庭和办公室很有可能直接通过光纤连接到大型中心局式交换机。
 3. 模块化多机箱路由器：人们对部署由一组通过高速网络互连的路由器组成的路由器集群越来越感兴
    趣。例如，许多网络接入点通过FDDI链路或千兆交换机连接路由器。路由器集群，有时也称为多机
    箱路由器，正变得越来越有吸引力，因为它们允许根据流量需求进行增量扩展，后面会对此进行解
    释。
核心路由器的典型使用寿命为18 - 24个月，之后流量的增加常常导致互联网服务提供商（ISP）淘汰旧
一代路由器，更换新的路由器。多机箱路由器可以将核心路由器的使用寿命延长到五年或更长时间，因
为它允许ISP从小规模开始，然后根据流量需求向集群中添加路由器。
21世纪初，瞻博网络（Juniper Networks）率先推出了其T系列路由器，该系列允许将多达16个单机箱路
由器（每个路由器最多有16个端口）通过交换结构组装成一个实际上的256端口路由器。这个多机箱系统
的核心是一个可扩展的256×256交换系统。同样在21世纪初，思科网络（Cisco Networks）推出了自己
的版本——CRS - 1路由器。
1.4.1 衡量交换机成本


                                                       22
在研究交换机扩展之前，了解交换机最重要的成本指标是很有帮助的。我们确定了两个关键成本指标：
CrossBar交换结构中的交叉点数和匹配算法的复杂度。它们分别对应于CrossBar交换硬件和软件的复杂
度。
  ● 指标#1：交叉点数：在电话交换的早期，交叉点是电磁开关，因此CrossBar交换结构的(N^{2})个交
    叉点是一项主要成本。即使在今天，对于规模为1000的非常大型的交换机来说，这仍然是一项主要
    成本。但是，由于交叉点可以看作只是晶体管，它们在超大规模集成电路（VLSI）芯片上占用的空
    间非常小。
对于电子交换机来说，真正的限制在于集成电路的引脚数量。例如，考虑到目前引脚数量的限制约为
1000个，其中很大一部分引脚必须用于其他因素，如电源和接地，即使是一个500×500交换机的单个位
片也不可能封装在一个芯片中。当然，可以将几个CrossBar交换输入复用在一个引脚上，但这会使每个
输入的速度降低到引脚I/O速度的一半。
因此，虽然CrossBar交换结构确实需要(N^{2})个交叉点（对于足够大的N，这确实很重要），但对于N
值高达200的情况，大部分交叉点复杂度都包含在芯片内部。因此，人们会在芯片内部放置尽可能大的可
实现的CrossBar交换结构，然后将这些芯片互连起来形成更大的交换机。因此，组合交换机的主要成本
是引脚成本和芯片之间的链路数量。由于这两者是相关的（大多数引脚用于输入和输出链路），引脚总
数是一个合理的成本衡量指标。更精确的成本衡量指标会考虑引脚的类型（背板、芯片、电路板等），
因为它们的成本不同。
限制构建大型单片CrossBar交换结构的其他因素包括总线上的电容负载、调度器复杂度以及机架空间和
功率问题。首先，如果试图使用256个输入和输出总线的CrossBar交换方法构建一个256×256的交换
机，负载可能会导致无法满足总线的速度要求。其次，请注意，一些集中式算法，如iSLIP，需要(N^{2})
位的调度状态，随着N的增大，扩展性不佳。
第三，许多路由器受可用性要求的限制，每个线路卡只能放置少数几个（通常是一个）端口。同样，由
于功率和其他原因，对可以放置在一个机架中的线路卡数量通常也有严格要求。因此，端口数量较多的
路由器可能会受到封装要求的限制，需要采用多机架、多机箱的解决方案，即由几个较小的交换结构连
接在一起形成一个更大的组合路由器。
  ● 指标#2：匹配计算的复杂度：在扩展交换机规模时需要考虑的另一个成本是匹配计算的计算复杂
    度。如前所述，对于一个具有N个输入/输出端口的交换机，计算一个良好的CrossBar调度（二分匹
    配）的总时间复杂度至少为(O(N))（在SERENA的情况下），最高可达(O(N^{2.5} log W))（在
    MWM的情况下），其中W是边（VOQ）的最大可能权重（长度）。因此，随着交换机规模N越来越
    大，以线路速率（例如每8纳秒处理一个数据包）为其底层（单片）CrossBar交换结构计算调度变
    得越来越困难。这种时间复杂度的挑战，加上前面提到的减少交叉点数的需求，自然促使研究人员
    和行业采用以下分治方法（P15）。
1.4.2 构建大型交换机的分治方法

                                                          23
分治方法是将大型交换机构建为较小交换机的互连网络。这样一来，每个较小交换机中的匹配计算时间
复杂度变得足够小，可以以线路速率处理，并且，正如我们接下来将详细阐述的，由于交叉点数的
(O(N^{2}))缩放特性，交叉点数也会显著减少。在以下小节中，我们将描述一种这样的分治解决方案，并
简要提及另一种在互联网泡沫时代（20世纪90年代末到21世纪初）实际应用于网络产品中的解决方案。
第一种解决方案受到最初为电话网络中的电路交换提出的Clos网络的启发并基于此。第二种解决方案改
编自用于并行处理的互连网络的经典解决方案。
这两种解决方案的原始思想解决的问题与我们目前定义的交换问题不同。在这两种原始应用中，即电话
网络中的电路交换和并行处理器的互连网络中，在任何时刻哪个 “输入端口” 应与哪个 “输出端口” 配对
是明确的，无需计算。在电话网络中，已建立的电话会话唯一地定义了输入和输出之间的匹配；在互连
网络中，任何时刻通信处理器的配对是由其上运行的并行算法决定的。在这两种解决方案中，计算问题
是，给定大型（虚拟）交换机中输入和输出之间的期望配对，确定实现这种匹配的小型（物理）交换机
的配置。
这个计算问题与我们目前提出的CrossBar调度问题显然不同，CrossBar调度问题恰恰是计算这样的配
对。然而，由于原始应用中存在更严格的约束，这个计算问题在严格意义上不一定更简单。例如，其中
一个约束是所有这些小型交换机都没有缓冲区，因此任何两个不同的输入 - 输出端口对所采用的两条路
径必须不相交（即不能共享中间节点或链路），或者等效地，大型交换机必须是非阻塞的。相比之下，
在分组交换问题中，假设输入端口有足够的缓冲区，并且不施加这样的约束。
1.4.3 适用于中型路由器的Clos网络
尽管目前在VLSI技术中对交叉点的关注度不高，但我们对路由器可扩展交换结构的研究从回顾历史上最
早提出的可扩展交换结构方案开始。Charles Clos在1955年首次提出他的想法，以降低电话交换机中机
电交换的成本。幸运的是，该设计还减少了连接多个较小交换机所需的组件和链路数量。因此，在当今
的环境下它仍然有用。具体来说，瞻博网络的T系列多机箱路由器产品似乎使用了Clos网络，该产品在47
年后的2002年推出。
基本的Clos网络采用简单的分治方法（P15），通过三个阶段的交换来减少交叉点数，如图16所示。第一
阶段将N个总输入分成每组n个输入的组，每个由n个输入组成的组通过一个小型（n×k）交换机连接到第
二阶段。因此，第一阶段有(N/n)个 “小型” 交换机。




                                                   24
图16 三级Clos网络
第二阶段由k个交换机组成，每个交换机都是一个(N / n×N / n)的交换机。每个第一阶段交换机的k个输
出都按顺序连接到所有k个第二阶段交换机。更准确地说，第一阶段交换机i的输出j连接到第二阶段交换
机j的输入i。第三阶段是第一阶段的镜像反转，第二阶段和第三阶段之间的互连也是第一阶段和第二阶段
之间互连的镜像反转。从输出端向左看中间阶段的视图与从输入端看中间阶段的视图相同。更准确地
说，第一阶段的每个(N/n)个输出都按顺序连接到第三阶段的输入。
如果一个交换机在输入和输出空闲时，总是可以使用空闲资源通过交换机建立连接，那么这个交换机就
被称为非阻塞交换机。因此，CrossBar交换机总是非阻塞的，因为它可以选择对应输入 - 输出对的交叉
点，该交叉点不会用于任何其他对。另一方面，在图16中，每个输入交换机到中间阶段只有k条连接，并
且每个中间阶段到第三阶段的任何特定交换机只有一条路径。因此，对于较小的k值，很容易阻塞新的连
接，因为可能不存在从输入I到中间阶段交换机的路径，而该中间阶段交换机又有空闲线路连接到输出
O。
 ● 电话网络中的Clos网络：克洛斯发现，如果(k ≥2 n - 1)，那么得到的克洛斯网络确实可以模拟
   CrossBar交换机（即是非阻塞的），同时将交叉点数减少到(5.6 N \sqrt{N})，而不是(N^{2})。对于
   较大的N，这可以节省大量成本。当然，为了实现这种交叉点数的减少，克洛斯网络增加了两个额外
   的交换阶段延迟，从而增加了延迟，但这通常是可以接受的。
从图17中很容易看出克洛斯定理的证明。如果一个此前空闲的输入i希望连接到一个空闲的输出o，那么考
虑i连接到的第一阶段交换机s。在s中最多可能有(n - 1)个其他输入处于忙碌状态（s是一个n×k交换
机）。这些(n - 1)条忙碌的输入链路最多可以连接到(n - 1)个中间阶段交换机。
同样，关注输出o，考虑o连接到的最后阶段交换机T。那么，T最多可能有(n - 1)个其他输出处于忙碌状
态，并且这些输出中的每一个最多可以通过(n - 1)个中间阶段交换机连接。由于s和T都连接到k个中间阶
段交换机，如果(k ≥2 n - 1)，那么总是可以找到一个中间阶段交换机M，它有一条空闲的输入链路连接
                                                             25
到S，并且有一条空闲的输出链路连接到T。由于S和T被假定为CrossBar交换机或其他非阻塞交换机，因
此总是可以将i连接到M的相应输入链路，并将M的相应输出链路连接到o。
如果(k = 2 n - 1)，并且n设置为其最优值(\sqrt{N / 2})，那么（图16中所有较小交换机的）交叉点总数
变为(5.6 N \sqrt{N})。例如，对于(N = 512)，这将交叉点数从CrossBar交换机的420万个减少到三级克
洛斯交换机的516,096个。更大的电话交换机，如1号电子交换系统（No. 1. ESS），可以处理65,000个
输入，它使用八级交换机来进一步减小交叉点规模。
  ● 减小中间阶段规模：另一方面，对于使用VLSI交换机的网络而言，重要的是交换机的总数以及互连
    交换机的链路数量。请记住，最大可能规模的交换机是在VLSI中制造的，并且其成本是固定的，与
    它们的交叉点规模无关。例如，瞻博网络通过在第一阶段连接16个16×16的T系列路由器，使用克洛
    斯网络构建了一个有效的256×256多机箱路由器。
对于一个完全配置的多机箱路由器，使用标准的克洛斯网络需要在第一阶段有16个路由器，在第三阶段
有16个路由器，在中间阶段有(k = 2×16 - 1 = 31)个交换机。显然，瞻博网络可以（并且确实）通过将(k
= n)来降低这种配置的成本。因此，瞻博网络的多机箱路由器在中间阶段只需要16个交换机。这正是我们
前面提到的使用(3m)个大小为(\sqrt{m}×\sqrt{m})的小型交换机构建一个大型(m×m)交换机的结构。




图17 证明k = 2n - 1的克洛斯网络是非阻塞的
当k从(2 n - 1)减少到(n)时，克洛斯网络会发生什么变化呢？如果(k = n)，克洛斯网络不再是非阻塞的。
相反，它变成了所谓的可重排非阻塞网络。换句话说，只要能够重新安排一些现有的输入和输出之间的
连接，以使用不同的中间阶段交换机，新的输入i就可以连接到o。附录A中描述了相关证明和可能的交换
算法。对理论不感兴趣的读者可以放心跳过这部分内容。
                                                              26
附录A中A.3.1节所有数学内容背后的核心要点如下：第一，(k = n)显然比(k = 2 n - 1)更经济，因为它将
中间层交换机的数量减少了一半。然而，虽然克洛斯网络是可重排非阻塞的，但目前用于交换机调度的
确定性边着色算法似乎相当复杂。第二，电话呼叫的匹配证明假设所有呼叫同时到达输入端口；当有新
呼叫到达时，可能需要重新安排现有呼叫以适应新的路由。
  ● 克洛斯网络和多机箱路由器：克洛斯网络用于电话网络时，必须是非阻塞的或至少是可重排非阻塞
    的；而用于多机箱路由器（即分组交换机）时，则没有这样的要求，这是因为这两种应用之间存在
    一个关键区别：电话交换机没有缓冲区，而（小型）分组交换机有。粗略地说，要求克洛斯网络是
    非阻塞的或可重排非阻塞的，等同于要求它在每个时隙内在中间交换机之间实现完美的负载平衡，
    这在没有缓冲的情况下是必要的。
然而，在有缓冲的情况下，这种负载平衡只需要在统计意义上是完美的，我们接下来将明确这个概念。
我们现在使用带有缓冲区的小型交换机的三级克洛斯网络，构建一个在统计意义上完美负载平衡的大型
交换机。在这个构建中，如前所述，我们将(k = n)，这样每个阶段都有n个交换机。由于(N = n^{2})，在
这种情况下，我们使用(3N)个大小为(\sqrt{N}×\sqrt{N})的交换机组成了一个(N×N)的交换机。例如，当
(N = 1024)时，我们仅使用96个32×32的交换机就构建了一个巨大的1024×1024交换机。我们将证明，
在这个三级克洛斯网络中，无需为这个巨大的1024×1024二分图（每次时隙）计算匹配。相反，每个小
型交换机只需为其32×32的二分图计算匹配，这在计算上要容易得多。
在这个三级克洛斯网络中，可以使用以下简单的随机策略实现统计意义上的完美负载平衡。在这个策略
中，第一阶段的每个交换机将其所有传入数据包统计均匀地分配到第二阶段的所有n个交换机。然后，在
第二阶段的每个交换机处，每个传入数据包都必须转发到第三阶段交换机，该交换机连接着数据包的目
的输出端口。这不会在第二阶段造成负载不平衡，因为可以证明，如果第一阶段的每个交换机都能完美
地实现统计负载平衡，那么在第二阶段的每个交换机处，发往其每个输出端口的流量在统计上是相等
的。出于类似的原因，第三阶段也不会出现负载不平衡。
那么，第一阶段的每个交换机如何才能完美地实现统计负载平衡呢？这可以通过以下简单直观的方案实
现，但有一个需要注意的地方，稍后会解释。我们仅描述第一级单个交换机的操作，因为其他交换机的
操作是相同的；在这个交换机中，我们仅考虑单个输入端口的操作，因为其他输入端口的操作也是相同
的。该方案是，对于每个传入数据包，这个输入端口简单地将其转发到一个输出端口，并且相应地转发
到第二级的一个交换机，这个输出端口和交换机是随机均匀选择的。
这个方案的问题在于，由于负载平衡是在数据包逐个处理的基础上（即在数据包级别）进行的，属于同
一TCP流的两个连续数据包可能会通过第二阶段的两个不同交换机，并经历不同的排队延迟。结果，它
们可能会从（同一个）第三阶段交换机出发，并无序地到达目的地。这种数据包无序到达可能会不必要
地降低TCP流的吞吐量，因为TCP客户端（两端）可能会错误地将其解释为网络拥塞导致的数据包丢
失，并相应地做出反应（例如，减小拥塞窗口的大小）。在互联网泡沫时代，这通常不是一个问题，因
为当时提供给TCP流的平均数据速率通常相当低（通常为几十到几百Kbps），这种数据包重新排序问题
不会使其进一步大幅降低。实际上，在21世纪初，瞻博网络似乎采用了类似的数据包级负载平衡方案，
尽管由于其文档（可能是有意）含糊不清，无法确定其实际做法。
                                                          27
近年来，三级克洛斯网络大多被用于构建大型数据中心网络（可以看作是一个巨大的交换机）
（Loukissas等人，2008；Cao等人，2013；Ghorbani等人，2017；Zhao等人，2019）。然而，如今
TCP流的数据速率要高得多，尤其是在数据中心网络中，这种数据包重新排序问题再也不能被忽视。因
此，最近所有针对三级克洛斯网络的基于数据包的负载平衡解决方案都必须消除或减轻这个数据包重新
排序问题。例如，在Ghorbani等人（2017）的研究中，巧妙地利用了终端主机上现代网卡内置的数据包
重新排序功能，将无序数据包的比例降至最低。
另一种负载平衡方法，即所谓的TCP哈希（Keslassy，2004），不存在数据包重新排序问题。在这种方
法中，属于同一TCP流的所有数据包必须通过克洛斯网络走相同的路径（即链路和交换机）。在第一级
交换机的每个输入端口，这可以通过对每个传入数据包的TCP流标识符（源和目的IP地址、源和目的端口
以及协议标识）进行哈希运算来实现，得到一个介于1到N之间的值，该值对应于这个数据包应该转发到
的交换机输出端口。然而，虽然TCP哈希消除了数据包重新排序问题，但它可能会导致严重的负载不平
衡。例如，大象流（流量很大的流）中的所有数据包都走相同的路径，会使沿途的链路和交换机拥塞。
因此，最近在一种称为安全随机交换（SRS）（Yang等人，2017b）的负载平衡交换（LBS）方案中，
TCP哈希方法通过两种安全机制得到了增强，以有效地缓解这种负载不平衡问题。
1.4.4 适用于大型路由器的贝内斯网络（Benes Networks）
就像1号电子交换系统（No. 1. ESS）电话交换机能够交换65,000条输入链路一样，特纳（Turner，
1997）和查尼等人（Chaney等人，1997）有力地论证了互联网（至少最终）应该由少数几个大型路由器
构建，而非多个较小的路由器。这样的拓扑结构可以减少较小路由器之间用于连接的冗余链路，从而降
低成本；还能缩短最坏情况下的端到端路径长度，减少延迟并提升用户响应时间。
本质上，克洛斯网络（Clos network）仅通过三个阶段，在交叉点复杂度方面大致具有(N\sqrt{N})的缩
放特性。这种权衡和一般算法经验（P15）表明，人们应该能够实现(NlogN)的交叉点复杂度，同时将交
换机的级数增加到(logN)。这样的交换网络在理论上、电话网络以及并行计算行业中已存在多年。其他替
代方案，如蝶形（Butterfly）、德尔塔（Delta）、榕树（Banyan）和超立方体（Hypercube）网络，都
是广为人知的竞争者。
虽然这个主题涉及的内容广泛，但本章仅关注德尔塔网络和贝内斯网络。类似的网络在许多实现中都有
应用。例如，华盛顿大学的千兆交换机（Chaney等人，1997）使用了贝内斯网络，它可以看作是两个德
尔塔网络的组合。附录A的A.4节概述了德尔塔网络与其他类似网络（通常差异较小）之间的区别。
在下面的内容中，我们仅在理论、电话网络和并行计算的历史背景下描述德尔塔网络和贝内斯网络，在
这些背景中，交换机没有或只有很少的缓冲区。我们将解释在这样的背景下，网络设计需要解决由异常
负载模式导致的拥塞问题。然而，在路由器的背景下，（小型）交换机有足够的缓冲区，这些缓冲区以
与克洛斯网络类似的方式进一步缓解了这些问题。
  ● 德尔塔网络（Delta Networks）：理解德尔塔网络最简单的方法是通过递归。假设有N个输入在左
    边，并且这个问题要简化为构建两个较小的（((N / 2))规模）德尔塔网络的问题。为了便于简化，假

                                                             28
  设第一阶段使用2×2交换机。一种简单的方案（图18）是检查每个输入希望与之通信的输出。如果输
  出在上半部分（输出的最高有效位（MSB）为0），那么该输入被路由到上半部分的(N/2)规模的德
  尔塔网络；如果输出在下半部分（即MSB = 1），则该输入被路由到下半部分的(N/2)规模的德尔塔
  网络。
为了节省第一阶段双输入交换机的数量，将输入分组为连续的对，每对共享一个双输入交换机，如图18
所示。因此，如果一对中的两个输入单元要发往不同的输出半部分，它们可以并行交换；否则，其中一
个将被交换，另一个则被丢弃或缓冲。当然，对于两个较小的（((N/2))规模）德尔塔网络，可以重复相
同的递归过程，将它们分解为一层2×2交换机，后面跟着四个(N/4)规模的交换机，依此类推。德尔塔网
络的完整扩展如图19的上半部分所示。请注意，图18中的递归结构在图19的第一阶段和第二阶段之间的
连接中可以清晰看到。




图18 通过将构建N输入德尔塔网络的问题简化为构建两个（N/2）输入德尔塔网络的问题，递归构建德
尔塔网络
因此，将问题简化为2×2交换机需要(logN)个阶段；由于每个阶段有(N/2)个交叉点，二进制德尔塔网络
的交叉点和链路复杂度为(NlogN)。显然，我们也可以通过在第一阶段使用d×d交换机，并将初始网络分
解为d个规模为(N/d)的德尔塔网络来构建德尔塔网络。这将阶段数减少到(log_dN)，链路复杂度减少到
(nlog_dN)。考虑到超大规模集成电路（VLSI）的成本，制造具有尽可能大d值的交换芯片以降低链路成
本更为经济。
德尔塔网络和它的许多近亲（见A.4节），如榕树网络和蝶形网络一样，具有一个很好的特性，称为自路
由特性。对于二进制德尔塔网络，人们可以通过在第i阶段沿着与输出(o=o_1, o_2, ..., o_s)（以二进制表
示）中(o_i)的值相对应的链路，找到从给定输入到给定输出的唯一路径。从图18中应该可以清楚地看到
这一点，在图中我们在第一阶段使用最高有效位，在第二阶段使用第二位，依此类推。对于(d≥2)，将输
出地址写为基数为d的数字，并以类似的方式遵循连续的数字。
在图18中可以直观地看到，德尔塔网络具有可逆性。可以通过以相同的方式跟踪输入的比特位，从输出
追溯到输入。因此，在图18中注意到，在从输出到输入的过程中，输入的倒数第二位选择两个连续的第

                                                         29
一阶段交换机之间的路径，最后一位选择输入。这种可逆性很重要，因为它允许使用德尔塔网络的镜像
反转版本（见图19的下半部分），该版本具有与原始德尔塔网络类似的特性。




图19 上半部分是德尔塔网络（图18），下半部分是镜像反转的德尔塔网络。上半部分用于分配负载，下
半部分用于路由
德尔塔网络的一个问题是拥塞。由于从每个输入到每个输出都有唯一的路径，德尔塔网络显然不是一个
置换网络。例如，如果每对连续的输入都希望向同一个输出半部分发送一个单元，那么只有一半的单元
可以进入第二阶段；如果这种情况重复，只有四分之一的单元可以进入第三阶段，依此类推。因此，对
于某些输出请求的组合，德尔塔网络的吞吐量可能会降低到一条链路的吞吐量，而不是N条链路的吞吐
量。
显然，要使德尔塔网络在任意输入请求的置换情况下都不那么容易出现拥塞，一种方法是在输入和输出
之间添加更多路径。推广克洛斯网络（图16）中的思想，可以构建一个贝内斯网络（图19），它由两个
（(logN)深度）的网络组成：左半部分是标准的德尔塔网络，右半部分是镜像反转的德尔塔网络。从输出
端向左看右半部分：注意从最后阶段到倒数第二阶段的连接与第一阶段和第二阶段之间的连接是相同
的。
人们也可以通过递归（P15）来可视化贝内斯网络，方法是扩展图18，添加一个由2×2交换机组成的第三
阶段，并以与第一阶段交换机连接到中间两个（((N/2))规模）网络相同的方式，将这些第三阶段的交换
机连接到中间的两个网络（图20）。观察到这个递归过程可以直接创建图19，而无需创建两个单独的德
尔塔网络。




                                               30
图20 递归构建贝内斯网络
观察图20中贝内斯网络的递归版本与图16中克洛斯网络的相似之处。可以利用这种相似性证明，贝内斯
网络可以以类似于我们对克洛斯网络可重排非阻塞特性的证明方式（见前面的方框内容），对任意输出
请求的置换进行完美路由。
在两次迭代的每一次中，像之前一样，首先在图20的第一阶段和最后阶段之间进行完美匹配，然后选择
两个中间交换机之一。然而，与克洛斯网络的证明不同，这里的算法必须在（((N/2))规模的）贝内斯网
络中递归地遵循相同的路由过程。或者，整个过程可以用边着色来表述。最终的结论是，在贝内斯网络
中可以完美地路由任意置换；然而，这样做相当复杂，并且不太可能在最短的数据包到达时间内以低成
本完成。
然而，回想前面提到的观点，对于克洛斯网络，一种随机策略在期望意义上比更复杂的确定性边着色方
案更有效。类似于在克洛斯网络中随机选择中间交换机，回到图19，人们可以在第一半部分为每个单元
随机选择一个目的地。然后可以在第二半部分使用反向德尔塔路由，将单元从随机的中间目的地路由到
实际的目的输出端口。使用随机中间目的地的想法可以追溯到瓦利安特（Valiant，1990），他首次在
（单副本）超立方体中使用这种方法进行路由。
与在克洛斯网络中使用随机选择中间交换机的情况一样，可以证明（在期望意义上），只要输入或输出
链路不拥塞，内部链路就不会拥塞。直观地说，负载分担部分将所有从任何输入发往任何输出链路的流
量均匀地分布在图19上半部分的所有N个输出链路上。在第二半部分，由于镜像结构，所有链路的流量又
会汇聚到目的输出链路。
例如，考虑图19中最后阶段第一个交换机的上行链路。一个重要的结论是，这条上行链路将承载发往输
出链路1的一半流量。这是因为这条上行链路承载了来自路由和复制网络（镜像反转的德尔塔网络）上半
部分输入节点的所有发往输出链路1的流量。并且，根据分配网络（德尔塔网络的上半部分）的负载分担
特性，这是发往输出链路1的一半流量。同样，可以论证这条上行链路承载了发往输出链路2的一半流
量。因此，如果输出链路1和2没有饱和，那么最后阶段交换机1的上行链路也不会饱和。对于第二半部分
的任何内部链路，都可以进行类似的论证。
                                                31
1.4.5 负载均衡交换（Load-Balanced Switching）
负载均衡交换（LBS，Chang等人，2002a，b；Ding等人，2014；Jaramillo等人，2008；Keslassy，
2004；Lin和Keslassy，2010；Yang等人，2017a，b）由Chang等人（2002a，b）首次提出，随后由
其他人（如Ding等人，2014；Jaramillo等人，2008；Keslassy，2004；Lin和Keslassy，2010；Yang
等人，2017a）进一步发展，这是另一种扩展到更大和更快交换机的方法。它通过消除前面提到的两个主
要扩展障碍之一来实现这一目标：在克洛斯网络中，每个小型交换机仍然需要进行二分匹配计算，而在
LBS中，这些计算被完全避免。




LB swtich
LBS架构基于瓦利安特负载均衡（Valiant load balancing，Leslie，1982）的思想，该思想早于任何IP
交换机或路由器的设计。它们依赖两个交换阶段来路由数据包。图21展示了一个通用的两阶段负载均衡
交换机的示意图。第一交换阶段将输入端口的第一阶段连接到中间端口的中心阶段，第二交换阶段将中
间端口的中心阶段连接到输出端口的最后阶段。每个交换阶段执行一个预定的周期性连接模式序列（因
此，无需进行二分匹配计算！），使得每个输入在每个交换阶段与每个输出的连接时间为(\frac{1}{N}) 。
LBS的一种典型连接模式是循环移位。例如，以下是一个易于理解和验证的标准循环移位连接模式。第
一交换结构执行一个周期性的 “递增” 匹配序列：在任何时隙t，每个输入端口i连接到中间端口((i + t)
\text{ mod } N) 。Chang等人（2002a）表明，当流量到达过程满足某些温和的平稳条件时，这个周期
性的匹配序列可以将传入流量非常均匀地分布到所有中间端口。另一方面，第二交换结构将执行一个周
期性的 “递减” 匹配序列：在任何时隙t，每个中间端口i连接到输出端口((i - t) \text{ mod } N) 。
尽管Chang等人（2002a）最初提出的基本负载均衡交换机能够实现100%的吞吐量，但它存在一个严重
的问题，即数据包的发送顺序可能会严重混乱。在基本负载均衡交换机中，输入端口的连续数据包在到
达时会被分散到所有N个中间端口。通过不同中间端口的数据包可能会遇到不同的排队延迟。因此，其中

                                                                    32
一些数据包可能会无序地到达它们的输出端口。这对互联网流量是有害的，因为广泛使用的传输控制协
议（TCP）会错误地将无序数据包视为拥塞和数据包丢失的指示。因此，已经提出了许多解决方案
（Ding等人，2014；Yang等人，2017a，b）来解决这个问题。
1.5 扩展到更快的链路速度
上一节重点讨论了交换机如何在规模上进行扩展。本节研究交换机如何在速度上进行扩展。现在，单个
光纤通道的速度可能在某个时候会趋于稳定。在20世纪90年代后期，许多专家表示，由于静态随机存取
存储器（SRAM）和光学的基本限制，单个光纤通道的速度将被限制在OC - 768（40Gbps）。那时，
光纤的容量将用于产生更多的单个通道（例如，使用多个波长），而不是更高速度的单个通道。使用更
多通道只会在增加端口数量方面影响交换，并且可以使用上一节中的技术来处理。
在过去的20年里，这个预测在很大程度上是正确的：在撰写本文时，OC - 768（40Gbps）仍然是广泛
使用的最高链路速度，尽管有100Gbps链路速度的（以太网）产品。然而，历史经验告诉我们，个别应
用确实有可能提高其速度需求，并且技术也可能令人惊讶地跟上步伐，使链路速度从目前的100Gbps提
高到，比如说，20年后的1Tbps。因此，寻找扩展交换机速度的技术是值得的。有三种常见的技术：位
切片、使用短链路和使用随机化内存共享。
1.5.1 使用位切片（bit slicing）实现更高速度的交换结构
应对链路速度增加的最简单方法是使用更快的时钟速率来运行交换电子设备。不幸的是，光速度呈指数
增长，而专用集成电路（ASIC）的时钟速率每年仅增长约10%。然而，根据摩尔定律，芯片上的晶体管
数量每18 - 24个月翻一番，而成本不会增加。因此，应对链路速度增加的最简单方法是使用并行性。
假设可以构建一个每个链路速度为S的CrossBar交换结构。那么，为了处理速度为(kS)（k为常数）的链
路，一种设计可以使用k个CrossBar “切片”。对于从一条链路传来的每组k位，每一位分别发送到每个
CrossBar切片。因此，每个切片看到的链路速度降低为(kS/k = S)，从而可以切实可行地实现。当然，
这意味着重组逻辑的速度也必须能够扩展。
如果以确定性方式将位分配到各个切片（即第一个单元的第1位发送到切片1，第2位发送到切片2，依此
类推），则重组逻辑可以简化，因为它知道下一位将出现在哪个切片上。然而，必须注意避免同步错
误。调度器可以对所有切片做出相同的决策，这使得调度器易于构建。
瞻博网络的T系列（Semeria和Gredler，2001）使用四个活动的交换结构平面（即切片）。它还使用第
五个平面作为热备用，以实现冗余。由于每个平面都使用请求 - 授权机制，如果在超时时间内没有收到
授权返回，则可以检测到平面故障。此时，只有在故障平面中传输的单元会丢失，故障平面将被替换进
行维护，备用平面将被切换投入使用。
到目前为止，虽然很少讨论，但冗余和容错对于大型交换机设计至关重要，因为风险更大。如果一个小
型的8端口路由器出现故障，只会影响少数用户。但是，一个大型的256×256端口的路由器必须几乎始终
正常工作，具备内部冗余以掩盖故障。这是因为在外部使用第二个相同的路由器作为冗余成本太高。大
                                                    33
多数互联网服务提供商（ISP）要求核心路由器符合网络设备构建系统（NEBS，Network Equipment
Building System，2002）标准。通常，大型路由器预计每年最多有五分钟的停机。
1.5.2 使用短链路实现更高速度的交换结构
到目前为止，我们一直忽略了互连网络中各阶段之间链路的物理长度这一特性。链路有多种形式，从芯
片之间的串行链路到背板走线，再到不同线路卡之间的电缆连接。直观地说，链路长度很重要，因为长
导线会增加延迟并降低比特率，除非使用更昂贵的信令技术（如光信令）来进行补偿。
观察Delta网络和Clos网络可以发现，这些网络在各阶段之间至少使用了一些长导线，其长度与 (O(N))
成正比。然而，也有一些互连网络可以使用均匀的短导线进行封装。这些就是所谓的低维网格网络。这
种网格网络在并行计算领域有着复杂的历史，曾被Cray和英特尔的超级计算机使用。
最简单的低维网格是一维环面（1D torus），它基本上是一条节点线，其中最后一个节点也与第一个节点
相连，形成一个逻辑环（图22）。二维环面（2D torus）基本上是一个二维节点网格，其中每行或每列
的最后一个节点也与同一行或列的第一个节点相连。三维环面（3D torus）是将相同的概念扩展到三维
网格。




图22 如何使用短导线在物理上封装一维环面
即使是逻辑上为环的一维环面，也似乎有一条长导线连接第一个和最后一个节点（图22）。然而，有一
种巧妙的方法可以将这条长线路的长度分摊到所有节点上，即利用一个简单的自由度（P13），将前半部
分节点布置在环的正向路径上（图22），后半部分节点布置在反向路径上。虽然从A到B的导线长度可能
翻倍，但这样就不存在长导线了。同样的想法可以扩展到二维和三维环面，通过在各行各列重复这个思
路来实现。
与蝶形（Butterfly）或Delta网络一样，一维环面的问题在于它会受到拥塞的影响，因为两个输入之间只
有两条路径。它还存在高延迟的问题，因为某些节点对之间的传输需要经过 (O(N / 2)) 跳。使用三维环
面可以缓解拥塞和延迟问题。例如，在一个8×8×8的三维环面中，平均每个消息可以（Dally，2002）在
90条长度为6跳的路径中进行选择。
                                                          34
Avici TSR路由器（Dally，2002）就是一个使用三维环面构建的路由器示例。它最多可以处理560个线
路卡，短导线的使用使其封装非常紧凑。一个260线路卡的配置可以通过仅使用跳线连接相邻背板来实现
无电缆封装。560线路卡版本则在两排机架之间使用一组短电缆（Dally，2002）。
除了使用短链路，三维网格还提供了大量的备用路径以实现容错，并且能够以最小的额外成本进行增量
升级。相比之下，一些互连网络往往需要以2的幂次方进行扩展。
1.5.3 利用随机化进行内存扩展
在目前我们所介绍的所有交换机中，数据包在拥塞期间都必须存储在缓冲区中。一般的经验法则是，路
由器需要具备一个往返时间（Roundtrip Time，RTT）的缓冲量，以便拥塞控制算法能够在不导致数据
包丢失的情况下降低传输速度。虽然使用更好的高层拥塞控制算法可能可以突破这个限制，但目前来
看，TCP和随机早期检测（random early detection，RED）的组合仍然需要这么多的缓冲量。以200毫
秒作为往返延迟的保守估计，一个两太比特的路由器必须具备0.4太比特的数据包缓冲区。因此，随着链
路速度的增加，如果在拥塞控制方面没有创新，内存需求也会相应增加。
考虑一个输入缓冲交换机，数据包以OC - 768的速度进入。因此，一个最小尺寸的数据包每8纳秒到达
一次，并且至少需要两次访问内存：第一次是存储数据包，第二次是将其读出以便通过交换结构进行传
输。在撰写本文时，最快的动态随机存取存储器（DRAM）的周期时间为50纳秒，显然，要使用DRAM
满足内存带宽需求，唯一的方法是使用更宽的内存字长。
不幸的是，内存字长不能大于最小尺寸数据包的长度，因为无法保证下一个数据包能同时被读出。也可
以使用静态随机存取存储器（SRAM，在撰写本文时其周期时间为4纳秒，理论上应该刚好够用），但这
样一来，成本会增加4到10倍。
解决这个困境的一种方法是使用并行的DRAM存储体。使用12个并行工作的DRAM存储体，每个存储体具
有40字节的访问宽度，就有可能跟上链路速度。从直观上看，这似乎是可行的。对于任何输入的数据包
流，可以将第一个数据包发送到DRAM 1，第二个发送到DRAM 2，依此类推。然而，由于服务质量
（QoS）和调度算法的原因，数据包的读出顺序并不明确。因此，在某些时间段内，可能所有数据包都
只从少数几个DRAM中读出，从而导致内存带宽竞争，最终造成数据包丢失。
这种内存竞争问题在计算机架构师使用交错内存时是很常见的。例如，如果一个数组按顺序分布在多个
内存存储体中，那么间隔特定步长的访问（例如列访问）可能都会命中同一个存储体。一种可能的巧妙
解决方法是借鉴CYDRA - 2无跨步敏感内存（Rau，1991）设计者的思路。他们的想法是对存储请求进
行伪随机交错，使得除了对同一字的访问之外，任何访问模式都不太可能导致热点问题。实际上，最近
有几篇论文（Lin等人，2009；Wang等人，2010；Zhao等人，2009）提出了实现这种无跨步敏感甚至
对抗性随机内存交错的技术。
在路由器环境中，不是将数据包1发送到DRAM 1，而是将每个数据包随机发送到一个选定的DRAM。当
然，与所有随机交错方案（见前面关于Clos和Benes网络的部分）一样，数据包的重组会变得更加复
杂，需要（在SRAM中？）保存状态以便对这些数据包重新排序。
                                                          35
1.6 结论
本章介绍了构建交换机的多种技术，从小型共享内存交换机，到思科GSR中使用的输入排队交换机，再
到瞻博网络T130和Avici TSR路由器中使用的更大、更具可扩展性的交换结构。
从根本上说，PIM和iSLIP的主要思想是认识到，通过使用虚拟输出队列（VOQ），可以（总共使用
(O(N^{2})) 位或每个端口 (O(N)) 位的通信复杂度）切实可行地传递所有期望的通信模式，从而避免队首
阻塞（HOL blocking）。这些方案进一步表明，使用随机化（PIM）或通过每个端口的轮询指针实现公
平性（iSLIP），可以在 (N log N) 时间内完成最大匹配。这里我们假设，使用如2.2.2节中所述的优先级
编码器等硬件，这种随机化和轮询指针机制可以在 (O(1)) 时间复杂度内实现。
虽然 (N log N) 是一个很大的数字，但通过让N个端口并行执行，可以将时间缩短到log N（在PIM中）或
一个较小的常数（在iSLIP中）。鉴于log N较小，即使是这样的延迟也可以通过流水线处理，使其在最短
的数据包时间内运行。基本的经验是，即使是像匹配这样看似复杂的算法，通过随机化和硬件并行性，
也可以在最短的数据包时间内运行。通过位切片技术可以进一步实现速度扩展。
与PIM和iSLIP相比，SW - QPS是一个令人惊叹的成果。它每个端口的时间复杂度严格为 (O(1)) ，且无
需任何硬件支持。其每个端口的通信复杂度也是 (O(1)) ，而PIM和iSLIP的通信复杂度为 (O(N)) 。然而，
它在整体吞吐量和延迟性能方面比iSLIP更优。它通过充分利用所有相关的算法技术实现了这一点：随机
化、硬件并行性以及极端流水线处理（滑动窗口技术可以说是极端流水线处理的一种体现）。
基于分治方法的算法技术也可以处理更大的端口数量。对交换实际成本的理解表明，即使是简单的三级
Clos交换机，对于端口数量高达256的情况也能很好地工作。然而，对于更大的交换机规模，结合了
（2logN）深度Delta网络的Benes网络更适合这项任务。这两种可扩展交换结构的主要问题都在于调
度。并且，与PIM一样，在这两种情况下，都是通过简单的随机化来简化复杂的确定性算法。在Clos和
Benes网络中，结构上的本质相似性使得可以先进行随机化负载平衡步骤，然后从随机的中间目的地进行
确定性路径选择。
类似的思想也用于通过选择随机的中间线路卡或将给定数据包（信元）随机发送到某个DRAM存储体来
减少内存需求。淘汰开关（knockout switch）使用随机化的2×2集中器树来实现k选N的公平性。因此，
随机化在交换机实现中是一个非常重要的概念。
有趣的是，本章中描述的几乎每一个新的交换机理念都催生了一家公司。例如，Kanakia在贝尔实验室从
事共享内存交换机的研究，然后离开创立了Torrent公司。瞻博网络的创立似乎源于Sindhu关于一种新交
换结构的想法，可能是基于通过随机中间线路卡进行分段的技术。McKeown在iSLIP取得成功后创立了
Abrizio公司。Turner、Parulkar和Cox创立了Growth Networks公司，将Turner的Benes交换机理念商业
化，该公司后来被思科收购。Dally将他在低维网格上实现无死锁路由的想法成功地从Cray计算机应用到
了Avici的TSR路由器中。



                                                                    36
37

\\section{解耦式交换机架构}

随着AI大模型训练规模的爆发式增长，传统机箱式（Chassis）交换机架构面临着前所未有的扩展性挑战。在万卡甚至十万卡级别的GPU集群中，传统交换机的物理限制和性能瓶颈日益凸显。解耦式交换机架构（Disaggregated Switch Architecture）应运而生，通过将交换功能分解为独立的硬件组件，实现了超越传统物理限制的扩展能力。

\subsection{传统机箱式交换机的局限性}

传统高端交换机采用一体化机箱设计，将线卡（Line Card）、交换网板（Fabric Card）、控制引擎和管理模块集成在一个物理机箱内。这种设计虽然便于管理和部署，但在超大规模AI网络场景下存在根本性局限：

\begin{enumerate}
    \item \textbf{端口密度限制}：受限于机箱的物理空间和散热能力，单台交换机的端口数量存在上限。即使采用最高密度的交换芯片，单机箱也难以满足万卡集群对端口数量的需求。
    \item \textbf{背板带宽瓶颈}：机箱内部线卡与交换网板之间通过背板或中置背板（Midplane）连接，背板的信号完整性限制了单链路速率和总交换容量。
    \item \textbf{单点故障风险}：控制引擎、交换网板等关键组件的故障会影响整台交换机的运行，尽管可以通过冗余设计缓解，但成本和复杂度显著增加。
    \item \textbf{扩展不灵活}：当集群规模增长时，需要更换整机或增加机箱数量，无法实现平滑的功能扩展。
\end{enumerate}

\subsection{DSF（Disaggregated Scheduled Fabric）架构}

DSF（Disaggregated Scheduled Fabric，解耦式调度交换架构）是Meta等公司针对AI训练网络提出的新一代交换架构。其核心创新在于将传统交换机的功能组件解耦为独立的网络设备，通过分布式协作实现超大规模交换能力。

\subsubsection{架构组成}

DSF架构由两类核心组件构成：

\textbf{接口节点（IN, Interface Node）}：也称为架顶解耦交换机（RDSW, Rack Disaggregated Switch）。IN是网络对外连接的边界，负责处理与外部网络的路由交互、数据包的封装/解封装以及流量的 ingress/egress 处理。每个IN作为独立的网络设备运行，具备完整的三层路由能力。

\textbf{交换节点（FN, Fabric Node）}：也称为交换网板解耦交换机（FDSW, Fabric Disaggregated Switch）。FN是内部交换的核心，专注于高速数据转发，无需支持复杂的三层路由功能。FN之间通过高速链路互联，形成无阻塞或低阻塞的交换核心。

从外部网络视角看，由多个IN和FN组成的DSF系统表现为一台逻辑上的统一交换机，其外部端口数量等于所有IN外部端口的总和。例如，由128个IN（每个IN提供32个400G端口）和64个FN组成的DSF系统，可被视为一台拥有4096个400G端口的巨型虚拟交换机。

\subsubsection{控制平面设计}

DSF的控制平面基于分布式网络操作系统实现。以Meta的FBOSS为例，它通过FBOSS State DataBase（FSDB）实现跨节点的实时状态同步。所有IN和FN运行统一的网络操作系统实例，通过控制平面协议协同工作：

\begin{itemize}
    \item IN之间维护全互联的BGP会话，交换IPv6邻居状态等路由信息
    \item FN之间通过专用控制协议协调转发状态
    \item 集中式控制器负责全局配置下发和状态监控
\end{itemize}

\subsection{接口节点与交换节点的分离}

IN与FN的分离是DSF架构的核心理念，这种分离带来了多方面的优势：

\textbf{功能解耦}：IN专注于复杂的路由处理、ACL、QoS策略等网络边缘功能；FN则专注于高速、低延迟的数据交换。这种专业化分工允许针对各自场景优化硬件设计。

\textbf{独立扩展}：根据集群规模，可以独立增加IN数量（扩展端口数量）或FN数量（扩展交换容量）。例如，当需要连接更多GPU服务器时，只需增加IN节点；当内部交换容量成为瓶颈时，则增加FN节点。

\textbf{异构硬件}：IN通常采用深缓冲区（Deep Buffer）芯片（如Broadcom Jericho3-AI）以支持复杂的拥塞管理；FN则可采用简化设计的芯片（如Broadcom Ramon3）以降低成本和功耗。

\textbf{故障隔离}：单个IN或FN的故障不会影响整个交换系统的运行。通过适当的冗余设计，可以实现组件级的热插拔和在线更换。

\subsection{包喷洒（Packet Spraying）技术}

包喷洒是DSF架构实现负载均衡的关键技术。传统ECMP（等价多路径）路由依赖流哈希（Flow Hashing）进行负载分担，但在AI训练场景下面临严峻挑战：

\begin{itemize}
    \item \textbf{大象流（Elephant Flows）}：AI工作负载产生持续时间长、流量大的数据流，容易在特定链路上造成拥塞
    \item \textbf{低熵流量}：RDMA通信的流数量有限，哈希结果可能导致多条流映射到同一条路径
    \item \textbf{负载不均衡}：上述因素导致 fabric 链路利用率出现显著偏差
\end{itemize}

包喷洒技术通过以下机制解决这些问题：

\textbf{信元化与多路径分发}：进入DSF的数据包在IN处被分割为固定大小的信元（Cell），这些信元被均匀地喷洒到所有可用的FN路径上。由于信元粒度远小于数据包，可以实现更细粒度的负载分担。

\textbf{信用机制}：DSF采用基于信用的拥塞控制算法。Ingress IN向Egress IN动态请求信用令牌（Credit Token），只有获得信用后才发送信元。这种机制确保了：
\begin{enumerate}
    \item 流量发送速率与接收端处理能力匹配
    \item 网络内部不会出现持续性拥塞
    \item 不同流量类别可以获得差异化的服务质量
\end{enumerate}

\textbf{虚拟输出队列（VOQ）}：每个IN为每个目标端口维护独立的VOQ。VOQ确保发往不同目的地的流量互不干扰，避免了队首阻塞（HOL Blocking）问题。每个VOQ独立调度，实现了细粒度的流量管理。

\textbf{顺序保持}：虽然信元通过不同路径传输，但DSF硬件在Egress IN处负责信元的重新排序，确保数据包按原始顺序交付给终端主机。这种乱序传输、有序交付的模式结合了负载均衡优势和应用兼容性。

\subsection{现代交换芯片技术}

解耦式交换机架构的实现离不开现代高性能交换芯片的支持。当前主流的两款芯片——Jericho3和Ramon3——在DSF架构中扮演着关键角色。

\subsubsection{Jericho3-AI：深缓冲区智能交换机芯片}

Broadcom的Jericho3-AI芯片专为AI训练网络设计，是DSF架构中IN节点的首选芯片：

\textbf{深缓冲区（Deep Buffer）设计}：Jericho3-AI集成了大容量片上缓冲区（数十MB级别），远超传统数据中心交换机的缓冲区大小。深缓冲区带来的优势包括：
\begin{itemize}
    \item 支持长距离链路传输（数公里甚至数十公里），补偿光纤传输延迟
    \item 为大象流提供充足的缓冲头room，避免丢包
    \item 支持复杂的无损拥塞控制协议（如PFC、ECN）的有效实施
\end{itemize}

\textbf{RoCEv2与RDMA支持}：Jericho3-AI针对RoCEv2（RDMA over Converged Ethernet v2）协议进行了深度优化：
\begin{itemize}
    \item 硬件支持PFC（Priority Flow Control），实现链路级的无损传输
    \item 支持ECN（Explicit Congestion Notification）标记，配合端网卡的拥塞控制算法
    \item 优化的数据包处理流水线，降低RDMA操作的延迟
\end{itemize}

\textbf{高基数端口}：单芯片可提供高达432个800G端口或等效的端口组合，支持灵活的端口配置以满足不同规模集群的需求。

\subsubsection{Ramon3：高性能交换网芯片}

Ramon3芯片在DSF架构中主要用于FN节点，专注于提供低成本、高效率的内部交换：

\textbf{简化功能集}：相比Jericho3，Ramon3去除了复杂的三层路由功能，专注于高速数据转发。这种功能简化带来了：
\begin{itemize}
    \item 更低的芯片成本和功耗
    \item 更紧凑的转发流水线，降低延迟
    \item 更高的端口密度
\end{itemize}

\textbf{信元交换优化}：Ramon3针对信元级交换进行了硬件优化，支持高效的信元路由和转发，与DSF的包喷洒机制完美配合。

\textbf{灵活拓扑支持}：支持多种内部互联拓扑，可以根据集群规模和性能需求灵活配置FN之间的连接方式。

\subsection{多级DSF架构}

在实际部署中，DSF架构可以根据规模需求构建为多级拓扑：

\textbf{L1区域}：由GPU机架、IN（RDSW）和FN（FDSW）组成的基本单元。一个L1区域可以互联数千个GPU端口，内部采用无阻塞拓扑。

\textbf{L2区域}：通过第二层Spine DSF交换机（SDSW）互联多个L1区域。SDSW使用与FDSW相同的硬件，但承担汇聚角色。一个L2区域可支持18K+的GPU规模。

\textbf{L3互联}：当需要跨建筑或跨区域互联时，通过Edge DSF交换机（EDSW）和Super-Spine层实现更大规模的互联。这一层通常采用一定的超配比（如4.5:1）以平衡成本与性能。

\subsubsection{输入均衡模式}

DSF架构中的创新特性之一是输入均衡模式（Input Balanced Mode）。该机制确保任何DSF设备的输入带宽不超过输出带宽：

当链路故障发生时，受影响的设备会向全网通告其降级的转发能力。其他设备收到这一信息后，会按比例减少向该设备发送的流量。这种分布式协调机制确保了即使在故障场景下，网络内部也不会出现拥塞。

\section{软件定义交换机与开源生态}

传统交换机采用封闭的软硬件一体化设计，网络操作系统（NOS）与特定硬件深度绑定。这种架构限制了创新速度——新功能的部署往往依赖于设备厂商的路线图，时间周期以月甚至年计。过去十年，以微软SONiC和Meta FBOSS为代表的开源交换机操作系统，正在重塑这一格局。

\\subsection{案例研究：SONiC与软硬件解耦}

微软通过SIGCOMM论文和Open Compute Project（OCP）推动了网络硬件与软件的彻底解耦。SONiC（云端开放网络软件）的出现，标志着交换机操作系统从封闭走向透明。

SONiC的核心创新在于基于容器化的微服务架构。传统的单体式NOS将所有的功能（路由协议、ACL、QoS等）紧密耦合在一起，任何小改动都可能影响整个系统。SONiC将这些功能拆分为独立的容器化服务，每个服务可以独立开发、测试和部署。这种架构允许开发人员在不影响转发平面的情况下更新路由协议或监控组件，新特性的发布周期从数月缩短到数周。

配合SAI（交换机抽象接口），SONiC可以运行在不同厂商的芯片之上。SAI定义了一套统一的API，将底层芯片的硬件细节抽象出来。这意味着，使用Broadcom芯片的交换机和使用Barefoot Tofino芯片的交换机，可以通过相同的SAI接口与SONiC交互。这极大地降低了云厂商对特定硬件供应商的依赖，并促进了网络技术的快速迭代。

\\subsection{案例研究：Meta FBOSS——交换机即服务器}

Meta的自研软件栈FBOSS采用了与SONiC不同的设计理念：将交换机视为一台特殊的"服务器"。FBOSS运行在一个精简的Linux发行版之上，利用Meta成熟的软件部署工具链进行管理。

这种设计带来了几个显著优势：
\\begin{enumerate}
    \\item \\textbf{运维一致性}：网络团队可以使用与服务器团队相同的工具（配置管理、监控、日志收集）来管理交换机，降低了运维复杂性。
    \\item \\textbf{快速迭代}：FBOSS的更新可以像服务器软件一样频繁进行。Meta能够在SIGCOMM 2018中展示的那样，快速部署新特性（如IPv6的全网支持），并能够应对流量在两年内增长30倍的严峻挑战。
    \\item \\textbf{深度可观测性}：由于运行标准Linux，FBOSS可以集成丰富的系统级监控工具，提供比传统NOS更细粒度的可见性。
\\end{enumerate}

\\subsection{案例研究：Meta Minipack——单芯片架构的极致追求}

在2019年的SIGCOMM研究中，Meta介绍了其F16数据中心架构。Meta认为，传统的由多块芯片构成的底座式交换机（Chassis Switch）不仅功耗高，且软件维护复杂。

为此，Meta推出了Minipack——一种基于单块高容量ASIC的128端口100G交换机。通过这种单芯片架构，Meta将机房内的网络跳数显著减少，并将每台设备的功耗和空间占用降低了50\\%。F16拓扑将整个机房划分为16个独立的平面，这种高度并行的设计确保了单点故障不会影响整体的连通性。

Minipack与FBOSS的结合，展示了软硬件协同设计的威力。由于只需要管理单一芯片，FBOSS的软件栈可以大幅简化；而FBOSS的灵活架构又使Minipack能够快速适应新的网络需求。这种"简单硬件+灵活软件"的模式，正在挑战传统网络设备厂商的复杂硬件策略。

\\subsection{趋势展望：可编程芯片与P4}

开源NOS的兴起与可编程交换芯片的普及相辅相成。Intel Tofino、Broadcom Trident3等芯片支持P4语言编程，允许用户自定义数据包处理流水线。SONiC和FBOSS都在积极拥抱这一趋势，通过集成P4编译器和运行时，使用户能够部署自定义的网络功能。

未来的交换机将越来越像服务器——运行开源操作系统、支持可编程硬件、通过标准工具管理。这一趋势不仅降低了网络创新的门槛，也为AI时代需要的高度定制化网络功能（如针对All-Reduce优化的拥塞控制）提供了实现基础。
% Google Cloud数据中心网络架构补充
% 第3章: 交换机设计原理与技术
\subsection{Google Cloud数据中心网络架构演进}

Google作为全球领先的云服务商，其数据中心网络架构经历了十余年的持续演进，从早期的传统Clos架构发展到今天基于光电路交换(OCS)的动态拓扑网络。本节将深入分析Google Jupiter网络的架构设计、关键技术及其对现代数据中心网络的启示。

\subsubsection{Jupiter网络架构概览}

Jupiter是Google构建的大规模数据中心网络架构，旨在为数万台服务器提供高带宽、低延迟的统一网络互联。根据Google在SIGCOMM 2022发表的论文\cite{jupiter2022}，Jupiter网络经历了从传统Clos拓扑到直接连接拓扑的重大架构演进。

\textbf{核心设计目标：}
\begin{itemize}
    \item \textbf{规模扩展性}：支持单个数据中心内数十万台服务器的互联
    \item \textbf{高带宽密度}：每服务器提供100+ Gbps的网络带宽
    \item \textbf{低延迟通信}：数据中心内任意两点延迟低于100微秒
    \item \textbf{增量部署}：支持不停机情况下的网络扩展和技术升级
    \item \textbf{成本优化}：在保持性能的同时降低资本支出和运营成本
\end{itemize}

\textbf{网络规模数据（截至2022年）：}
\begin{itemize}
    \item 支持超过6 Pb/秒的数据中心总带宽
    \item 覆盖数十万台服务器的统一网络fabric
    \item 10年间实现5倍速度和容量提升
    \item 资本支出降低30\%，功耗降低41\%
\end{itemize}

\subsubsection{从Clos到OCS：拓扑架构演进}

传统的大规模数据中心网络通常采用三层Clos拓扑结构：机架顶部交换机（ToR）→汇聚层（Aggregation）→骨干层（Spine）。然而，这种架构面临以下挑战：

\begin{enumerate}
    \item \textbf{技术升级困难}：当网络从40Gbps升级到100Gbps时，预部署的Spine交换机成为瓶颈，限制了新服务器的性能发挥
    \item \textbf{资源利用率不均}：固定的拓扑结构无法适应动态变化的流量模式
    \item \textbf{扩展成本高昂}：每次扩展都需要预先购买和部署最大规模的Spine层设备
\end{enumerate}

\textbf{Google的创新解决方案——光电路交换(OCS)：}

Google开发了基于MEMS（微机电系统）的光电路交换技术，将Jupiter从Clos拓扑演进为直接连接拓扑：

\begin{itemize}
    \item \textbf{消除Spine层}：移除了传统的骨干交换层，直接在汇聚块之间建立动态连接
    \item \textbf{OCS核心组件}：使用可旋转的MEMS镜子动态映射光输入端口到输出端口
    \item \textbf{动态拓扑重构}：通过软件控制实现网络拓扑的灵活调整，无需物理重新布线
\end{itemize}

\begin{table}[htbp]
\centering
\caption{Clos vs OCS架构对比}
\begin{tabular}{|l|c|c|}
\hline
\textbf{特性} & \textbf{传统Clos} & \textbf{OCS直接连接} \\
\hline
拓扑结构 & 固定三层 & 动态可重构 \\
扩展方式 & 预部署Spine & 增量添加汇聚块 \\
重配置时间 & 数周（物理布线） & 毫秒级（软件控制） \\
平均路径长度 & 3-4跳 & 1.4跳 \\
直接路径流量占比 & - & 60-80\% \\
功耗 & 基准 & 降低41\% \\
\hline
\end{tabular}
\end{table}

\subsubsection{Jupiter的光电路交换技术}

\textbf{OCS工作原理：}

光电路交换通过物理调整光路来实现端口间的连接，无需光电转换：

\begin{enumerate}
    \item \textbf{输入光信号}：来自源交换机的光信号进入OCS设备
    \item \textbf{MEMS镜子阵列}：两组可在二维平面旋转的MEMS镜子控制光路走向
    \item \textbf{光路重构}：镜子角度调整将光信号反射到指定的输出端口
    \item \textbf{无包处理}：OCS仅反射光信号，不解析包头或进行路由决策
\end{enumerate}

\textbf{技术优势：}
\begin{itemize}
    \item \textbf{超低延迟}：光信号直接通过，引入的延迟可忽略不计（纳秒级）
    \item \textbf{高带宽}：单根光纤可支持多波长传输（WDM），总带宽达数Tbps
    \item \textbf{低功耗}：无需光电转换和包处理，功耗远低于电交换机
    \item \textbf{可扩展性}：8个OCS机架可支持任意逻辑拓扑的数据中心网络
\end{itemize}

\textbf{拓扑工程与流量工程协同：}

Google采用集中式SDN控制器对Jupiter进行统一管控：

\begin{itemize}
    \item \textbf{拓扑工程(TE)}：根据预测的流量模式调整OCS连接，优化网络拓扑
    \item \textbf{流量工程(TM)}：在固定拓扑内通过路由算法优化流量分布
    \item \textbf{协同优化}：结合拓扑和流量工程，在直接连接架构上实现接近Clos的吞吐量
\end{itemize}

\subsubsection{Google Titanium子系统与Axion处理器}

\textbf{Titanium基础设施处理单元(IPU)}

2024年，Google推出了Titanium子系统，这是一种面向云数据中心的基础设施卸载架构，与AWS Nitro系统理念相似但实现方式有所不同：

\begin{table}[htbp]
\centering
\caption{Titanium IPU核心功能}
\begin{tabular}{|l|l|l|}
\hline
\textbf{功能域} & \textbf{卸载内容} & \textbf{性能指标} \\
\hline
网络卸载 & 虚拟交换、安全组、负载均衡 & 支持100Gbps+网络 \\
存储卸载 & 块存储I/O、加密、校验 & 650K IOPS \\
安全卸载 & 系统启动信任根、加密 & 硬件级安全 \\
管理卸载 & Hypervisor、Borg调度器 & 释放主CPU资源 \\
\hline
\end{tabular}
\end{table}

\textbf{Axion ARM处理器}

Google Axion是Google首款自研的数据中心ARM处理器，于2024年Google Cloud Next大会发布：

\begin{itemize}
    \item \textbf{架构}：基于Arm Neoverse V2，支持Armv9指令集
    \item \textbf{性能}：
    \begin{itemize}
        \item 比当前最快的通用Arm实例快30\%
        \item 比可比的x86实例性能好50\%
        \item 能效比x86高60\%
    \end{itemize}
    \item \textbf{集成}：与Titanium IPU深度集成，主CPU专注于客户工作负载
    \item \textbf{实例类型}：C4A（计算优化）、N4A（成本优化，2025年GA）
\end{itemize}

\textbf{与AWS Nitro的对比：}

\begin{table}[htbp]
\centering
\caption{Google Titanium vs AWS Nitro对比}
\begin{tabular}{|l|c|c|}
\hline
\textbf{特性} & \textbf{Google Titanium} & \textbf{AWS Nitro} \\
\hline
发布年份 & 2024 & 2017 \\
处理单元 & IPU (与Intel合作) & Nitro Cards (自研+合作伙伴) \\
CPU架构 & Axion (ARM) & Graviton (ARM) + x86 \\
网络性能 & 100Gbps (C4A) & 200Gbps (最新代) \\
存储IOPS & 650K & 数百万 (最新代) \\
安全芯片 & Titan (信任根) & Nitro Security Chip \\
\hline
\end{tabular}
\end{table}

\subsubsection{对现代交换机设计的启示}

Google Jupiter和Titanium的演进为现代数据中心网络设计提供了重要启示：

\begin{enumerate}
    \item \textbf{软硬件协同设计}：网络架构需要与上层应用（如AI训练、大数据分析）协同优化，而非孤立设计
    
    \item \textbf{动态拓扑的价值}：OCS技术证明，通过软件定义的光交换可以实现物理拓扑的灵活调整，同时保持高性能
    
    \item \textbf{分层卸载架构}：Titanium展示了如何将基础设施任务（网络、存储、安全）从主CPU卸载到专用硬件，提升整体效率
    
    \item \textbf{ARM在数据中心的崛起}：Axion处理器标志着Google正式加入ARM服务器阵营，与AWS Graviton、Microsoft Cobalt等形成竞争格局
    
    \item \textbf{AI原生网络设计}：随着AI工作负载占比增加，网络设计需要特别考虑大象流、低熵流量等AI训练特有的流量模式
\end{enumerate}

% ============================================================
%  现代网络设备的演进：从ASIC到可编程芯片
% ============================================================
\section{现代网络设备的演进：从ASIC到可编程芯片}

以太网从1970年代诞生至今，历经半个多世纪的变革，但一件有趣的事始终没有改变：构建
网络平台的基石——交换机——仍旧是今天每一个数据中心最不可或缺的设备。与此同时，这
块"基石"本身正在经历一场安静却深刻的革命：从封闭的专有硬件，走向开放的白盒生态；
从固化的ASIC，走向可用代码重新定义数据面的可编程芯片。理解这条演进轨迹，是理解
现代云数据中心网络架构的前提。

\subsection{交换机的本质：从集线器到智能转发}

要理解为什么交换机如此重要，我们先回到它出现之前的世界。

早期以太网使用的核心设备是\textbf{集线器（Hub）}。集线器工作在OSI模型的第一层
（物理层），本质上是一个信号放大器：一个端口收到信号，立刻广播到所有其他端口。这
意味着，如果有$n$台机器接在同一个集线器上，它们共享同一条碰撞域（collision domain），
任意时刻只能有一台机器发送数据。当$n$变大时，碰撞概率急剧上升，CSMA/CD机制虽然
能检测和重传，但有效吞吐量会骤降。这是一个根本性的扩展瓶颈。

交换机（Switch）的出现解决了这个问题。交换机工作在\textbf{数据链路层（第二层）}，
核心能力是基于\textbf{MAC地址}进行精准转发。它内部维护一张\textbf{MAC地址表}
（也叫转发表、CAM表），记录"某个MAC地址在哪个端口方向"。当一帧到达时，交换机
查表，找到目的MAC对应的端口，只把这帧送往那一个端口，而不是广播出去。

这个"查表-精准转发"的设计带来了质的改变：
\begin{itemize}
  \item 每对通信端口之间拥有\textbf{独享带宽}，互不干扰；
  \item 将广播域限制在必要范围内，有效\textbf{隔离广播风暴}；
  \item 支持\textbf{全双工通信}，理论带宽翻倍；
  \item 可以在同一时刻处理\textbf{多对端口的并发转发}，背部总线带宽决定了整机吞吐上限。
\end{itemize}

以一个直观的数字说明差距：以24口100\,Mbps集线器为例，当碰撞概率超过50\%时，有效吞
吐仅约为理论带宽的30\%（即2.4\,Gbps总带宽中实际可用约720\,Mbps）；同规格的二层
交换机则每个端口独享100\,Mbps，24口可同时全速工作，有效带宽接近100\%的线速利用率。
这就是为什么交换机取代集线器之后，局域网带宽"感觉"提升了好几倍——不是信道变快了，
而是碰撞消失了。

MAC地址表的学习过程非常直观：当某端口收到一帧，交换机"记住"源MAC地址和对应端口；
若目的MAC在表中已有记录，则精准转发；若没有，则暂时泛洪到所有端口，等待目的方的回应。
这个自学习机制使得交换机无需人工配置即可正常工作，极大简化了部署。

\subsubsection{交换技术的演进：从共享到独占}

交换机内部的交换方式经历了三代演进，每一代都针对前一代的瓶颈做出改进。理解这条演进
脉络，有助于我们理解为什么今天的数据中心交换机会做出某些特定的设计选择。

\textbf{第一代：端口交换（Port Switching）}是集线器时代的小修小补。它将所有端口分成
若干组，同组内端口共享一条总线，不同组之间通过一个中央交换矩阵互联。碰撞域从一个大
域被切成了若干小域，但本质上仍是共享介质——同组内的端口竞争带宽的问题依旧存在。这
一代方案解决了规模问题，却没有解决效率问题。

\textbf{第二代：帧交换（Frame Switching）}是真正意义上的"每端口独享带宽"方案，也是
现代以太网交换机的核心模式。一旦交换机知道目的端口，就可以建立一条专属的逻辑信道，
不再与其他端口竞争。然而，"什么时候开始转发"这个问题，在不同场景下有三种不同答案：

\begin{itemize}
  \item \textbf{存储转发（Store-and-Forward）}：等收完整帧、验过FCS校验和再转发。好处是
    可以过滤掉错帧，不把错包传播到下一段链路；代价是引入了与帧长度成正比的排队延迟。
    对于一个1500字节的以太帧，在1\,Gbps链路上这个延迟约为$12\,\mu\text{s}$，在10\,Gbps
    链路上则降至约$1.2\,\mu\text{s}$；
  \item \textbf{直通转发（Cut-Through）}：只读取前6字节的目的MAC地址，立刻开始向目的端
    口转发——此时帧的其余内容还在源端口接收中。延迟极低，接近$1\,\mu\text{s}$甚至更低，
    非常适合对延迟敏感的高频交易等场景；但错帧、碰撞帧会被忠实地转发出去；
  \item \textbf{无碎片转发（Fragment-Free）}：读取前64字节再转发。64字节是以太网最小帧
    长，也是CSMA/CD能够可靠检测碰撞的最短帧。因此这种方式能过滤掉碰撞产生的"碎片帧"，
    同时比存储转发的等待时间短很多——是二者的工程折中。
\end{itemize}

\textbf{第三代：信元交换（Cell Switching）}来自ATM（异步传输模式）世界。它把可变长度
的帧切成固定长度的信元（ATM标准规定53字节：5字节头 + 48字节载荷），好处是调度硬件
设计极大简化——每个信元大小相同，缓冲区和交换矩阵的设计可以用简单的定时器驱动。代价
是切割和重组开销不可忽视，53字节中5字节是头部，开销比接近10\%。随着以太网速率从
百兆提升到万兆，ASIC足以应对可变帧长的存储转发，信元交换的优势消失，ATM在局域网
领域逐渐退出历史舞台。

今天的高性能数据中心交换机几乎都采用\textbf{存储转发}模式。原因在于：现代ASIC的流水
线处理能力已经如此之强，一个64字节最小包的存储转发延迟已经压缩到$300\sim500\,\text{ns}$，
与直通转发的延迟差距已经可以忽略；而CRC校验带来的错帧过滤能力，对保障数据中心的链路
质量至关重要——尤其在运行RDMA的无损网络中，一个错帧引发的重传可能触发严重的性能抖动。

\subsubsection{三层交换机：打通二三层的桥梁}

传统路由器工作在第三层（网络层），负责IP包的路由转发。但如果在同一栋楼里不同VLAN
之间通信，每次都要把流量送到楼外的路由器再转回来，这明显是浪费。

三层交换机的出现解决了这个问题：它在二层交换能力之上，集成了IP路由的硬件查找功能。
第一个数据包走软件路径完成路由决策，建立转发条目；后续同流数据包直接由ASIC硬件以
线速完成"一次路由，多次交换"（Route Once, Switch Many）。这种\textbf{硬件路由}
能力使得VLAN间通信的延迟和吞吐与二层交换相差无几，成为数据中心接入层和汇聚层的
标配方案。

\subsection{路由器的架构：控制面与数据面的分离}

路由器与交换机最大的不同，不在于端口速度或价格，而在于它承担了一项更复杂的使命：
在\textbf{不同网络之间}做出最优路径决策。为了理解这一点，我们先看路由器内部的
功能分层。

现代路由器在逻辑上分为两个平面：

\textbf{控制平面（Control Plane）}：负责运行路由协议，构建和维护路由表。典型协议包括：
\begin{itemize}
  \item \textbf{OSPF（Open Shortest Path First）}：链路状态协议，域内路由，采用
    Dijkstra最短路径算法，收敛速度快；
  \item \textbf{BGP（Border Gateway Protocol）}：路径向量协议，域间路由，支撑着整个
    互联网的路由体系；AS之间通过BGP交换可达性信息，每个AS可以实施
    丰富的路由策略（Route Policy）；
  \item \textbf{ISIS（Intermediate System to Intermediate System）}：同样是链路状态
    协议，在运营商骨干网和大型数据中心中广泛使用，扩展性强于OSPF。
\end{itemize}

\textbf{数据平面（Data Plane）}：负责按照控制平面建立的转发表（FIB，Forwarding
Information Base）对每个到来的数据包做快速的IP地址匹配和转发动作。数据平面
的性能（转发速率、延迟）由专用硬件决定，通常是网络处理器（NPU）或
专用ASIC。

这种分离是有深刻原因的。路由协议的计算往往复杂：BGP要处理数十万条路由，OSPF要
运行SPF算法，这些适合在通用CPU上灵活执行。但一旦转发表建好，对每个数据包的查表
动作需要在纳秒级完成。我们来估算一下这个压力有多大：

\begin{quote}
1\,Tbps的线速转发，若数据包平均大小为\textbf{128字节}（加上20字节以太帧头部，
共约148字节），则每秒包数（PPS）约为：
\[
\text{PPS} = \frac{1 \times 10^{12}\,\text{bps}}{148 \times 8\,\text{bits/packet}}
           \approx \frac{10^{12}}{1184} \approx \mathbf{845 \times 10^6\,\text{pkt/s}}
\]
即\textbf{约8.5亿包/秒}。若进一步考虑最坏情况（64字节最小帧），则PPS将达到约
1.5\,\text{Gpps}。这意味着每个数据包分配到的处理时间不足\textbf{1纳秒}。
\end{quote}

在如此极端的时间约束下，即使是最简单的软件哈希查表操作（每次约\(10^2 \sim 10^3\)纳秒）
也望尘莫及。只有将查表逻辑实现为ASIC里的\textbf{硬件流水线}——在芯片时钟的每个
周期内完成一个包的头部解析和表项匹配——才能满足这个要求。这是专用ASIC不可取代
的根本原因。

\subsubsection{路由器内部结构}

一台高端路由器通常由以下核心模块组成：

\textbf{线卡（Line Card）}是路由器与外部网络的接口。每块线卡上集成了端口PHY（物理层
收发器）、MAC/PCS芯片以及一块转发ASIC（或NPU）。当数据包到达时，线卡的转发ASIC
负责完成三项工作：首先对报文做头部解析和完整性校验；其次在本地FIB（由路由引擎下发）
中查找目的IP的下一跳和出端口；最后打上内部标签，经交换网送往目的线卡。线卡的处理能力
决定了整机的"前端速度"，典型高端线卡可支持\textbf{2\,Tbps}的接口总带宽。

\textbf{交换网（Switch Fabric）}是连接所有线卡的内部高速互联网络，相当于路由器的
"内部骨干"。传统设计采用共享总线或共享内存，所有线卡公用同一块交换资源，当线卡数量
增多时，交换网容量很快成为瓶颈。现代高端路由器普遍采用\textbf{多级Clos结构}的交换网
——这正是下文要介绍的胖树（Fat-Tree）拓扑思想的内部应用。思科CRS-1路由器（2004年
发布）内部采用三级Clos交换网，最高支持92个线卡槽位、总交换容量达到92\,Tbps，其内部
结构本质上就是一个按机框级别构建的分布式Clos网络。

\textbf{路由引擎（Route Engine）}是路由器的"大脑"，运行在通用CPU（通常是x86或ARM）
上。路由引擎维护两张表：\textbf{RIB（路由信息表，Routing Information Base）}收录所有
路由协议学到的路由条目，是"原始数据库"；\textbf{FIB（转发信息表，Forwarding
Information Base）}是从RIB中优选出的最优路由，以可供线卡ASIC直接查询的格式存储，
是"转发决策表"。路由引擎将FIB实时同步到各线卡——这个下发过程通常通过内部总线
（如HA总线）完成，速度比外部网络快得多。路由引擎故障（如崩溃重启）\textbf{不影响
已建立的转发状态}：线卡本地的FIB副本继续运作，已有流量照常转发——这就是控制面与
数据面分离的最直接价值体现。

\subsubsection{从路由器到SDN：控制面集中化的思想根源}

传统路由器的控制面"分布式"运行在每台设备上，靠OSPF/BGP等协议在网络设备间协同。
这种方式鲁棒性好，但也带来了问题：每台设备都需要理解复杂协议，配置繁琐，网络行为
难以整体编排。

2006年，斯坦福大学Martin Casado和Nick McKeown提出了\textbf{OpenFlow}的概念：
将控制平面从每台设备中提取出来，集中到一个中央控制器（Controller）；转发设备
只负责执行控制器下发的\textbf{流表（Flow Table）}，实现"数据面哑化，控制面智能化"。
这正是SDN（Software Defined Networking，软件定义网络）思想的核心。

\subsection{芯片的战争：Broadcom ASIC与可编程芯片的路线之争}

如果说交换机是数据中心网络的基石，那么交换芯片（Switch ASIC）就是这块基石的心脏。
理解芯片市场的格局，是理解现代交换机价格、性能、可编程性的关键。

\subsubsection{交换芯片的市场格局}

今天，市面上主流的数据中心交换机——无论是思科、华为、H3C的品牌交换机，还是
广达（Quanta）、Celestica等ODM厂商的白盒交换机——其核心交换芯片几乎都来自
同一家公司：\textbf{Broadcom（博通）}。

Broadcom的Trident系列（接入/汇聚层）和Tomahawk系列（核心层）占据了数据中心交换
芯片市场超过80\%的份额。这并非偶然：博通在交换芯片领域积累了数十年的工艺和IP，其
芯片的功能完整性、稳定性和单位带宽成本都形成了显著的领先优势。

排在其后的竞争者包括：
\begin{itemize}
  \item \textbf{Marvell（美满）}：主打低功耗，在运营商以太网和接入交换机领域有
    较强存在感；其Prestera系列在中低端市场有稳定的用户群；
  \item \textbf{Intel（英特尔）}：通过收购Barefoot Networks进入可编程交换芯片领域
    （后文详述）；
  \item \textbf{国产芯片}：包括以太链路智能、华为海思等，近年来随着信创需求快速发展，
    但在超大规模数据中心的生产部署中仍在追赶阶段。
\end{itemize}

\subsubsection{ASIC的性能边界与"协议无关"的梦想}

Broadcom类的交换ASIC是典型的\textbf{协议相关（Protocol-Dependent）}架构：芯片出厂
时，转发流水线的每个阶段已经被固化——解析以太头、VLAN、IP头、MPLS标签，查ARP表、
FIB表，执行QoS标记……这些处理步骤以硬件电路的形式刻在硅片上，无法修改。

这种设计的优势是性能极致：12.8\,Tbps的单芯片带宽，256个400GbE端口，转发延迟低至
$300\sim500\,\text{ns}$，同时功耗控制在合理范围内。但代价也很明显：\textbf{一旦出现
新协议、新功能需求，就必须等待下一代芯片}。芯片从立项到量产通常需要3-5年，这在
云时代已经是难以接受的周期。

举一个具体的例子。VXLAN协议2011年才提出标准，但早期的Broadcom Trident系列芯片
并不能硬件卸载VXLAN的封装/解封装。阿里、腾讯等互联网公司在部署VPC网络时，不得不
在服务器上用软件（Open vSwitch或DPDK）来做VXLAN处理，付出了大量CPU资源。直到
后续版本的Trident才陆续加入硬件VTEP支持。

这个案例揭示了一个本质矛盾：\textbf{云和互联网业务的创新速度远快于芯片硬件的迭代
周期}。这正是可编程交换芯片诞生的动力。

\subsubsection{可编程芯片：FPGA和P4的崛起}

可编程网络处理的尝试并不新鲜。\textbf{FPGA（Field Programmable Gate Array）}长期
被用于网络设备的原型验证和小批量定制场景——它可以用硬件描述语言（VHDL/Verilog）
重新定义电路逻辑，实现任意协议。但FPGA的单位带宽成本远高于ASIC，且编程模型对网络
工程师不友好，无法大规模部署。

\textbf{NPU（Network Processing Unit）}是另一条路线——专为包处理设计的多核微处理器，
如EZchip、Netronome等。NPU可编程性好，但受限于软件执行的本质，在超高速包处理场景
（如400G/800G）上与ASIC相比存在明显的性能差距。

真正让"可编程高速交换芯片"成为现实的，是Barefoot Networks于2016年推出的
\textbf{Tofino芯片}，以及与之绑定的\textbf{P4语言}。

\subsection{P4与Barefoot TNA：可编程数据平面的工程实践}

2013年，斯坦福大学和普林斯顿大学的研究人员提出了P4语言（Programming Protocol-independent
Packet Processors）。P4的核心思想是：网络工程师应该能够\textbf{用高层语言描述数据平
面的行为}——如何解析报文头、如何匹配-动作，而不必深入到RTL级别的硬件描述。

P4语言脱胎于一种叫做\textbf{PISA（Protocol-Independent Switch Architecture）}的
抽象模型：

\begin{center}
\begin{tabular}{|l|l|l|}
\hline
\textbf{阶段} & \textbf{功能} & \textbf{P4对应构造} \\
\hline
报文解析（Parser） & 从比特流中提取头部字段 & \texttt{parser} 声明 \\
入口Match-Action & 对头部字段做匹配，执行动作 & \texttt{control} + \texttt{table} \\
流量管理（TM） & 排队、调度、ECN标记 & 硬件固化，P4不直接控制 \\
出口Match-Action & 出口方向的二次处理 & \texttt{control} + \texttt{table} \\
反解析（Deparser） & 将修改后的头部重新打包成帧 & \texttt{control deparser} \\
\hline
\end{tabular}
\end{center}

这个抽象模型的精妙之处在于：它和具体的硬件\textbf{解耦}了。同一份P4代码，理论上可以
编译到Barefoot Tofino芯片、也可以编译到FPGA或软件模拟器上运行。

为了让读者对P4语言有直观感受，下面给出一个最小化的\texttt{parser}声明示例，展示如何
解析以太网帧和IPv4包头：

\begin{verbatim}
/* 定义以太网头部结构 */
header ethernet_t {
    bit<48> dstAddr;          // 目的MAC地址
    bit<48> srcAddr;          // 源MAC地址
    bit<16> etherType;        // 上层协议类型（0x0800=IPv4）
}

header ipv4_t {
    bit<4>  version;
    bit<4>  ihl;
    bit<8>  diffserv;
    bit<16> totalLen;
    bit<8>  ttl;
    bit<8>  protocol;
    bit<32> srcAddr;
    bit<32> dstAddr;
}

/* Parser：定义报文解析的状态机 */
parser MyParser(packet_in pkt,
                out headers hdr,
                inout metadata meta) {
    state start {
        pkt.extract(hdr.ethernet);          // 从报文中提取以太网头
        transition select(hdr.ethernet.etherType) {
            0x0800 : parse_ipv4;            // etherType=0x0800跳转IPv4解析
            default : accept;               // 其他协议直接接受（不解析）
        }
    }
    state parse_ipv4 {
        pkt.extract(hdr.ipv4);              // 继续提取IPv4头
        transition accept;                  // 解析完成，进入Match-Action阶段
    }
}
\end{verbatim}

这段代码对应PISA流水线的第一个阶段（Parser）。可以看到P4的核心特征：
\begin{itemize}
  \item \textbf{状态机驱动}：Parser由若干\texttt{state}组成，通过\texttt{transition select}
    按字段值决定跳转——这里\texttt{etherType}字段决定是否进入IPv4解析状态；
  \item \textbf{显式提取}：\texttt{pkt.extract(hdr.X)}明确指定从报文比特流中提取哪个
    头部结构，填入对应的PHV（Packet Header Vector）寄存器，供后续Match-Action阶段使用；
  \item \textbf{协议无关}：若要新增对VXLAN的支持，只需增加一个\texttt{state parse\_vxlan}
    并在\texttt{etherType}/\texttt{protocol}判断中加入对应跳转——无需等待芯片厂商。
\end{itemize}

\subsubsection{Tofino原生架构（TNA）}

Barefoot Networks（后被Intel收购）于2016年6月推出第一代Tofino交换芯片。Tofino不是
简单地在ASIC上加一个CPU，而是从头设计了一个\textbf{完全可编程的转发流水线}。

TNA（Tofino Native Architecture）的核心结构如下：

\textbf{Pipeline与端口的映射}：一块Tofino芯片包含多个独立的Pipeline（流水线）。
每个物理端口属于固定的一个Pipeline，不可跨Pipeline共享处理资源。以典型的4-Pipeline
Tofino为例：1-16端口属于Pipeline 0，17-32端口属于Pipeline 1，以此类推。这意味着
\textbf{在网络设计中，需要将同一子卡的端口规划在同一Pipeline内}，否则跨Pipeline
流量必须经过回环（Recirculate）处理，占用额外带宽。

\textbf{流水线内的阶段（Stage）}：每个Pipeline包含12个可编程的Match-Action单元（MAU，
Match-Action Unit），ingress和egress流各走一遍，共24个逻辑阶段，但硬件资源是共享的。
每个Stage包含固定配额的：
\begin{itemize}
  \item \textbf{SRAM}：用于精确匹配表（Exact Match），如VXLAN VNI到VRF的映射；
  \item \textbf{TCAM}：用于通配符匹配（Wildcard Match），如最长前缀匹配（LPM）路由查找；
  \item \textbf{PHV（Packet Header Vector）}：存放从报文中提取的头部字段的寄存器组，
    是Match-Action操作的基本载体。
\end{itemize}

这些资源是\textbf{全局有限的}，P4程序员的核心工作之一就是在有限的Stage资源中合理
规划表项大小，避免某个Stage成为瓶颈。

\subsubsection{阿里云Tofino实践：SNA高性能云网关}

2017年，阿里云开始基于Tofino研发第一代高性能云网关SNA（Smart Network Appliance），
2019年正式上线。这是国内互联网公司将可编程交换芯片规模落地的最早工程实践之一。

为什么阿里云的网关业务会选择PISA架构的Tofino，而不是更成熟的Broadcom ASIC？答案
在于业务的特殊性：

云网关需要同时处理\textbf{大促模式}和\textbf{公有云模式}两种差异极大的流量场景。
大促场景下，流表规模相对有限，但需要极高的PPS（每秒包数）吞吐；公有云模式下，
租户数量巨大，流表规模需求极高。固化ASIC只能做有限的妥协，而Tofino可以通过重新
编译P4程序，在两种模式下分别优化表项分配——这正是为什么SNA维护了两个版本的P4
程序（大促版本和公有云版本）。

一个典型的资源权衡问题：SNA云网关使用的是4-Pipeline Tofino芯片，但实际对外提供的
带宽只有3.2\,Tbps而非4×1.28\,Tbps=5.12\,Tbps。这是因为其中一个Pipeline被专门用于
Recirculate（报文回环处理），以支持需要多次查表才能完成的复杂业务逻辑。这个"用
带宽换可编程灵活性"的取舍，在固化ASIC上是不可能实现的。

\subsubsection{Tofino vs Broadcom Trident4：硬件资源对比}

为了让读者对两种架构有更直观的认识，这里做一个简要对比：

\begin{center}
\begin{tabular}{|l|l|l|}
\hline
\textbf{特性} & \textbf{Barefoot Tofino} & \textbf{Broadcom Trident4（TD4）} \\
\hline
架构类型 & 可编程PISA & 协议固化ASIC（Tile架构）\\
流水线 & 12 Stage/Pipeline，可编程 & 13 ingress + 3 egress Stage，固化 \\
TCAM用途 & 任意自定义匹配字段 & 固化用于路由/ACL查找 \\
协议扩展 & P4重编译即可 & 需等待下一代芯片 \\
单芯片带宽 & 6.4\,Tbps（Tofino2） & 12.8\,Tbps \\
包转发速率（PPS） & $\sim$2.0\,Gpps（Tofino2） & $\sim$3.2\,Gpps（TD4） \\
典型功耗 & $\sim$300\,W & $\sim$320\,W \\
性能/功耗比 & 略低（可编程开销） & 极致优化（固化流水线） \\
典型应用 & 网关、DPI、自定义协议 & 标准二/三层交换 \\
\hline
\end{tabular}
\end{center}

这两条技术路线并非你死我活的竞争关系，而是在不同场景下互补：Broadcom ASIC用于
标准转发的大规模部署，Tofino系列用于需要快速迭代自定义协议逻辑的场景。

至此，可编程芯片（P4/Tofino）解决了数据平面的灵活性问题——我们终于可以用代码重新定义
转发流水线，而不必等待下一代芯片。然而，数据中心网络还有另一个根本性的问题悬而未决：
\textbf{谁来控制这些设备？} 每台交换机各自独立运行OSPF/BGP的分布式模式，真的是最好的
答案吗？如果可以把控制权集中到软件层，是不是能做更聪明的调度？这个问题将网络界分成了
泾渭分明的两派，引发了软件SDN与硬件SDN之间旷日持久的路线之争。

\subsection{软件SDN与硬件SDN：路线之争的本质}

SDN的故事从OpenFlow开始，但OpenFlow很快发现自己陷入了一个两难困境。理解这个
困境，需要先区分两种截然不同的SDN实现路线。

\subsubsection{硬件SDN：以交换机为VTEP的方案}

\textbf{硬件SDN}是数据中心SDN市场的主流方案，主要由思科、华为、H3C等网络设备厂
商推动。其核心做法是：将VXLAN的VTEP（VLAN Tunnel End Point，即VXLAN封装/解封
装点）部署在物理交换机上，服务器只负责发原始流量，由交换机完成overlay封装。

这种方案的优势很明显：转发性能强（交换机ASIC硬件处理VXLAN封装），服务器CPU
零开销。但随着云规模扩大，硬件SDN的问题逐一暴露出来，可以归纳为三个层面的矛盾。

第一是\textbf{规模与迭代速度的矛盾}。交换机ASIC内置的TCAM和SRAM容量有限，通常
支持数万到数十万条流表项；而大型公有云的VM数量动辄以百万计，每个VM的ARP/MAC
状态都需要占用一条表项。硬件SDN网关无法支撑这个规模。更麻烦的是，一旦需要新增功能
（如支持新的隧道协议），就必须等待芯片厂商的下一代产品——而这个等待通常以年计，与
互联网公司按周迭代的节奏完全脱节。

第二是\textbf{可靠性的矛盾}。在硬件VTEP方案中，VXLAN封装集中在少数几台核心SDN
网关设备上。这些设备一旦出现故障，整个POD的overlay网络都会受影响。传统网络设备
的冷升级（割接窗口）往往需要提前安排业务停机，这对"7×24小时不中断"的云服务来说
是难以接受的成本。

第三是\textbf{软硬件深度耦合}。商业SDN控制器与下层交换机硬件之间的接口往往是私有的，
采购一家厂商的控制器就意味着深度绑定该厂商的交换机平台。这种"锁定效应"不仅拉高了
硬件采购成本，也使得运营商难以引入竞争来降低价格——这与互联网公司追求的"白盒化"
思路背道而驰。

值得注意的是，2020年中国SDN市场规模已达25.9亿元，预计2025年达到120亿元，年均
增速超过30\%。市场在高速增长，但增长的主力是\textbf{软件SDN和混合方案}，而非纯粹
的硬件SDN——大量用户在经历了上述问题后，正在逐步迁移到以服务器端虚拟交换机为核心
的软件定义架构。

\subsubsection{软件SDN：把VTEP搬到服务器上}

软件SDN的核心思路与硬件SDN相反：将VTEP部署在\textbf{服务器或智能网卡}上，
利用Open vSwitch（OVS）、DPDK-vSwitch或eBPF等技术在服务器内核或用户态完成
VXLAN的封装和解封装。

这种方案的优势是灵活性极强：
\begin{itemize}
  \item 租户网络完全用软件定义，表项规模理论上无上限（受服务器内存约束）；
  \item 新协议支持只需升级软件，可以做到不停机滚动更新；
  \item 每台服务器都是独立的VTEP，天然分布式，没有中心单点；
  \item 硬件选型自由——物理交换机只需做标准IP转发，无需特殊SDN功能。
\end{itemize}

阿里云的VPC网络（第三章已详细介绍）就是软件SDN路线的典型代表：HAVS（硬件加速
虚拟交换）将OVS卸载到专用硬件，服务器CPU几乎零开销，同时保留了软件定义的全部灵活性。

\subsubsection{OpenFlow的局限与混合方案的现实}

OpenFlow是SDN最初的标准接口协议：控制器通过OpenFlow协议向交换机下发流表，流表
规定了对特定报文执行的匹配和动作。这个思路理论上完美，但工程实现中遇到了现实障碍：

\textbf{流表规模问题}。OpenFlow v1.0的流表项匹配字段长达252位，是普通IPv4路由TCAM
匹配字段（60-80位）的3倍以上。支持完整OpenFlow流表的TCAM成本极高，现有交换机
芯片能支持的OpenFlow流表条目非常有限（通常只有数千到数万条）。

OpenFlow v1.1引入多级流表（Multi-Table Pipeline）来缓解这个问题：将一条大流表项
分解为多个小表项，利用表间流水线处理降低总条目数。举例来说，若有1万种流需要按
``源IP/目的IP''和``端口/协议''两个维度共同匹配，合并为一张全字段表需要1万条；若
分解为一级``源IP/目的IP''表（1000条）和二级``端口/协议''表（1000条），通过流水线
串联组合即可覆盖相同的$1000 \times 1000 = 1,000,000$种匹配空间，而总条目仅需
2000条。这利用了乘法组合效应：$(M+N) \ll M \times N$，两级各1000条的组合能力
远超单级1万条。

然而即便如此，传统网络设备厂商对OpenFlow的态度始终暧昧：一方面声称支持，另一方面
刻意模糊"软件定义网络"和"网络设备可编程"的概念边界，不愿让控制权真正转移到第三
方控制器。这种局面使得真正意义上的OpenFlow大规模部署，在2015年之后逐渐沉寂。

\textbf{Google B4：迄今最成功的SDN商业应用。}2012年，Google公开宣布在其内部WAN
上实现了流量调度系统B4。B4使用Google自研的定制化分布式交换机，结合SDN和OpenFlow
实现了\textbf{End Application Flow Control}，带宽利用率接近100\%，是传统商业路由器
方案的2-3倍。B4分为Site层和Global层，Site间运行BGP/ISIS，Global层的TE（流量工程）
控制器统一调度跨站点带宽，优先级感知地分配流量——这是OpenFlow在广域网流量工程领域
迄今为止最令人印象深刻的工程实践。

\textbf{路线之争的真正答案：混合而非对立。}回顾硬件SDN和软件SDN的争论，超大规模互联
网公司最终走向了一条\textbf{分层混合}的路线，而非非此即彼的选择。其逻辑是：\textbf{物理
转发层}（underlay）用白盒ASIC交换机做高性能L2/L3转发，追求线速和低延迟；\textbf{租户
虚拟网络层}（overlay）则用软件SDN——将VXLAN封装/解封装部署在服务器端的OVS、DPDK
或智能网卡上，实现无限扩展的流表规模和零停机的快速迭代。两者分工明确：交换机只负责
高速转发IP包，不感知任何租户逻辑；租户逻辑全部封装在服务器的SDN软件层。这种架构既
保留了ASIC的极致性能，又获得了软件SDN的无限灵活性——后文介绍腾讯自研交换机时，读者
会看到这一混合路线在工程中的具体落地。

\subsection{腾讯自研交换机：互联网公司的硬件自主之路}

2014年，腾讯发表了一篇内部文章，宣布推出自研交换机Tswitch。这个时间节点很有意思——
彼时，Amazon、Google、Alibaba等互联网巨头已经悄然开始了自己的交换机自研历程，整个
行业的"自研浪潮"正在形成。

\subsubsection{为什么要自研交换机？}

在回答"怎么做"之前，我们先理解"为什么"。

传统数据中心交换机由Cisco、H3C、华为等OEM厂商提供，软硬件全部封闭。这对快速迭代
的互联网公司而言，存在几个根本性的痛点：

\textbf{功能迭代慢，无法快速响应业务需求。}商业交换机厂商面向全球客户维护庞大的功能
集，一个版本的开发周期通常以年计。互联网公司的网络需求日新月异（虚拟化、VPC、SDN、
大规模BGP等），等待厂商支持往往来不及。以OSPF为例，白皮书可能有150页，但腾讯实际
用到的子特性不超过10个——为了这10个特性，还要购买支持了150页特性的昂贵设备。

\textbf{成本高昂，性价比低。}互联网公司购买商业交换机，可能只用了不到10\%的功能。
自研方案通过购买廉价ODM白盒裸机（广达、Celestica、Foxconn等厂商），在上面运行自
主研发的网络OS，以极低的成本进入万兆时代。百度就以这种方式，相比千兆商业交换机，
以更经济的方式完成了万兆数据中心的建设。

\textbf{故障诊断能力弱，严重依赖厂商。}商业交换机出现问题，查根因依赖厂商支持，
响应周期长。自研交换机拥有完整的源码和调试能力，能快速定位并修复网络故障。

\subsubsection{白盒交换机的生态链}

理解腾讯自研交换机，需要先理解整个白盒交换机产业生态的分层结构：

\textbf{芯片层}：Broadcom、Marvell提供交换ASIC，SDK（软件开发工具包）免费开源给
客户，包含硬件表项的增删改查接口（如路由表、下一跳表的操作）。

\textbf{硬件层（ODM）}：广达（Quanta）、天弘（Celestica）、富士康（Foxconn）、
智邦（Accton）等ODM厂商生产标准化白盒交换机裸机（Bare Metal Box）。这些设备在
硬件规格上与OEM品牌交换机差距不大——因为大家用的是同样的Broadcom芯片。

\textbf{软件层}：互联网公司的核心竞争力所在。软件方案分为三个子层次：
\begin{enumerate}
  \item \textbf{网络系统汇聚层（硬件适配层）}：对芯片SDK进行封装，向上暴露统一的
    业务接口，屏蔽不同芯片厂商的SDK差异。换芯片厂商时，只需修改这一层；
  \item \textbf{协议层}：OSPF、BGP、STP等二三层协议的实现；
  \item \textbf{应用管理层}：CLI、SNMP、gRPC等管理接口。
\end{enumerate}

当时主流的白盒NOS（网络操作系统）授权方案如下：

\begin{center}
\begin{tabular}{|l|l|l|}
\hline
\textbf{软件方案} & \textbf{使用公司} & \textbf{备注} \\
\hline
ICOS & Amazon, Alibaba, Microsoft & 功能最完整，有授权费用 \\
Zebos & Tencent & 有授权费用 \\
Pica8 & Baidu & 有授权费用 \\
Cumulus & Facebook & 有授权费用 \\
Quagga & 院校科研机构 & 免费开源，仅L3路由 \\
\hline
\end{tabular}
\end{center}

腾讯选择了Zebos作为初期的NOS基础，并在此之上构建Tswitch的上层特性和运维能力。
Zebos的选择有其工程逻辑：相比ICOS，Zebos的授权门槛较低，且允许深度定制修改源码，
符合腾讯"快速起步，深度自研"的策略；相比完全开源的Quagga，Zebos提供了工业化程度
更高的L2/L3完整协议栈（含STP/MSTP/OSPF/BGP），可以直接支撑数据中心生产部署，无需
从零补齐二层协议。

\subsubsection{腾讯自研交换机的技术实践}

腾讯Tswitch的演进路线展示了互联网公司自研交换机从"能用"到"好用"再到"智能化"
的成熟路径。

\textbf{转发平面的持续增强。}腾讯Tswitch的功能迭代路线，清晰地反映了数据中心网络在
过去十年的技术演进轨迹。

早期的NOS 1.0以满足基础二三层功能为目标：二层需要STP/MSTP、LACP链路聚合、802.1Q
VLAN；三层需要OSPF/BGP路由协议。这是运行一个生产数据中心的最低门槛。随着业务规模
的扩张，Tswitch逐步引入了几类重要的现代数据中心协议：

其一是\textbf{无损网络支持（RoCE/RDMA）}。AI训练、高性能存储等业务对网络延迟和
可靠性要求极高，要求交换机支持PFC（Priority Flow Control，基于优先级的流量控制）和
ECN（Explicit Congestion Notification，显式拥塞通知）。PFC在检测到端口队列深度超过
阈值时，向上游发送PAUSE帧，暂停该优先级流量的发送，从而实现零丢包传输；ECN则在
队列占用率达到水位时，主动在报文IP头打标（CE位），让发送端提前降速，避免队列满溢。
两者联合使用，可以将大规模GPU训练集群的网络丢包率控制在\textbf{百万分之一}以下。

其二是\textbf{EVPN/VXLAN硬件卸载}。在VPC网络场景下，服务器端的SDN软件完成
VXLAN封装，物理交换机只需做标准IP转发；但在部分场景（如裸金属服务器、容器集群），
交换机需要硬件卸载VXLAN封装/解封装。Tswitch通过升级Broadcom SDK驱动，实现了
在Trident系列ASIC上的VTEP功能，单芯片支持超过\textbf{16K个VNI（VXLAN Network
Identifier）}，满足中等规模租户隔离需求。

其三是\textbf{gRPC控制通道}。传统交换机管理依赖SSH + SNMP/CLI，配置变更需要逐台
登录执行命令。Tswitch引入gRPC作为统一的南向接口，控制器可以批量推送配置和查询状态，
实现数千台交换机的秒级并发配置下发——这是实现Primus集中式路由方案的基础设施前提。

\textbf{可视化与立体监控。}传统商业交换机对运营人员近乎"黑盒"——交换机内部到底发生
了什么，往往只能通过SNMP轮询（每5分钟采一次数据）或者事后查日志来推断。在一个
拥有数千台交换机、数十万条链路的数据中心里，这种"事后诸葛亮"的排障方式意味着：从
故障发生到定位根因，平均耗时以\textbf{小时}计，严重时甚至以天计。

腾讯在Tswitch中构建了一套多层次的网络可观测体系：

\textbf{Telemetry（遥测）}以订阅/推送模型替代了轮询SNMP。交换机每\textbf{1秒}（而非
传统的5分钟）向采集平台推送一次端口速率、队列深度、CPU/内存利用率等状态，使得网络
状态的时间分辨率提升了整整300倍。当一条链路在某次推送窗口内出现异常，告警系统几乎
实时响应，将故障发现时间从分钟级压缩到\textbf{秒级}。

\textbf{IFA（In-situ Flow Analysis）}解决了一个更棘手的问题：端到端延迟的逐跳归因。
传统做法是在源端/目的端各打时间戳，测量往返延迟，但无法判断延迟发生在哪一跳。IFA
在每个中间交换机上，将该跳的入队时间戳和出队时间戳嵌入报文本身（利用报文头的
保留字段），使得接收方可以还原完整的逐跳延迟分解。在一次大促压测中，腾讯工程师
借助IFA数据发现某条核心链路的单跳排队延迟从正常的$30\,\mu\text{s}$飙升到
$3\,\text{ms}$，精准定位到一块存在微突发问题的交换机端口。

\textbf{MOD（Micro-burst On-Demand）}专门应对"微突发"（Micro-burst）这一数据中心
网络的隐形杀手。微突发是指持续时间在\textbf{毫秒以内}的瞬时流量尖峰，持续时间短到
SNMP轮询和常规Telemetry都无法捕获，但却足以将队列瞬间填满、触发丢包。MOD在交换机
本地维护一个高精度（微秒级）的环形缓冲区记录队列深度变化，当检测到队列深度超过
阈值时，立即导出该时间窗口的原始数据，供后端分析系统做事后根因分析。

\textbf{Netsense}通过主动探测补充了被动监测的盲区。Netsense向网络中定期注入探测包
（类似ICMP，但使用专用协议以精确控制发送时间），检测交换机间路径可达性和往返延迟。
相比等待业务报文触发的被动告警，Netsense能在\textbf{无流量区间}提前发现链路异常，
将链路故障感知时间缩短到\textbf{1秒以内}。

这四个工具联合使用，使得腾讯网络运维团队对数据中心网络的"感知颗粒度"从粗放的5分钟/
POD级别，提升为精细的秒级/端口级别——这正是"白盒化运营"的核心价值：拥有源码，
才能拥有这种级别的可观测性。

\textbf{集中式路由方案（Primus）。}传统BGP的收敛速度依赖协议本身的定时器机制，
在大规模拓扑下往往较慢。腾讯自研了Primus集中式路由系统：在规整的叶脊拓扑中，
中央路由控制器直接感知链路状态，指导驱动层快速更新转发表。实测数据显示，
\textbf{直连故障收敛时间提升2-3倍，跨跳故障收敛时间提升近10倍}，从分钟级降至秒内。

Primus的工作原理如下：在规整的叶脊（Leaf-Spine）拓扑中，每台交换机通过gRPC将
实时链路状态持续上报给Primus控制器；控制器在收到链路变更事件后立即执行SPF（最短
路径优先）计算，生成差量转发表（Delta FIB）；最后通过gRPC将更新直接推送到受影响
交换机的驱动层，驱动层调用Broadcom SDK接口直接更新ASIC硬件转发表，\textbf{完全
绕过了传统BGP的Hold Timer（默认90秒）、Keepalive超时和路由重计算的多轮等待}，
这是收敛速度大幅提升的核心原因。

\textbf{零接触配置（ZTP）。}在数据中心批量部署场景下，人工逐台配置交换机效率极低。
Tswitch实现了Zero Touch Provisioning：新设备上架后，开机自动通过DHCP获取管理IP，
连接ZTP服务器，自动下载配置和软件版本，全程无需人工干预。在一次有1000台设备的
新建POD中，ZTP将配置时间从数天压缩到数小时。

\textbf{热升级（Warm Reboot）。}传统交换机升级需要冷重启，业务中断10-30秒甚至更长。
Tswitch实现了基于Docker容器的热升级机制：系统层（Host OS）和应用进程（各协议Docker）
可以独立升级，在升级窗口内实现\textbf{零丢包}或极少量丢包，极大降低了维护成本。

热升级期间转发面保持不中断的关键在于两个机制：其一，\textbf{转发状态保持（Forwarding
State Preservation，FSP）}——转发面进程（ASIC SDK）在升级期间持续运行，路由协议
Docker重启后通过读取共享内存中保存的FIB快照重建路由状态，无需重新计算路由表；其二，
\textbf{优雅重启（Graceful Restart，GR）}——OSPF/BGP进程通过GR机制通知邻居设备进入
辅助模式（Helper Mode），邻居继续正常转发流量但暂停路由更新，直到本机协议进程重新
稳定并发出"重启完成"信号后，邻居才恢复正常的路由更新流程。两个机制协同工作，使得
协议Docker升级对网络流量几乎完全透明。

\subsubsection{自研交换机的价值总结}

回顾腾讯、阿里、Google、Facebook等互联网巨头的自研交换机实践，可以归纳出三条
共同的核心价值：

\begin{enumerate}
  \item \textbf{快速响应业务需求}：自研团队与业务团队同处一个组织，需求到特性的
    周期可以压缩到周级别，而不是传统厂商的年级别；

  \item \textbf{深度可观测性}：拥有完整源码意味着可以在任何关键路径加入监控和调试
    能力，实现"白盒化运营"，故障定位速度数量级提升；

  \item \textbf{降低成本，打破垄断}：白盒化后，硬件从垄断的品牌交换机变为竞争充分
    的ODM市场，价格大幅下降；同时，网络设备的核心价值从硬件转向软件，这正是
    互联网公司的优势所在。
\end{enumerate}

互联网数据中心的未来已经清晰：Broadcom/Marvell等半导体厂商提供成熟的ASIC芯片；
Quanta/Celestica等ODM厂商提供白盒硬件；互联网公司在上面运行自研的NOS和控制平面
软件。传统网络巨头的封闭时代，正在被这场"白盒+软件定义"的革命所终结。

\begin{tcolorbox}[colback=blue!5!white, colframe=blue!50!black,
  title={\textbf{核心概念速览：从ASIC到P4的技术脉络}}]
\begin{itemize}
  \item \textbf{ASIC交换芯片}：Broadcom Trident/Tomahawk，协议固化，性能极致，
    但不可编程；
  \item \textbf{P4语言}：描述数据平面行为的领域专用语言，目标是协议无关（Protocol-Independent）；
  \item \textbf{PISA架构}：Parser → Match-Action Pipeline → Deparser，可编程的抽象
    转发模型；
  \item \textbf{Tofino/TNA}：Barefoot（Intel）的可编程交换芯片，首个商用的P4靶标；
  \item \textbf{OpenFlow/SDN}：控制面与数据面分离，中央控制器统一管控，Google B4是
    最成功的实践；
  \item \textbf{白盒交换机}：ODM硬件 + 开源/自研NOS，互联网公司打破Cisco/华为垄断的
    主要路径。
\end{itemize}
\end{tcolorbox}

\begin{thebibliography}{99}
\bibitem{jupiter2022} Poutievski et al., ``Jupiter Evolving: Transforming Google's Datacenter Network via Optical Circuit Switches and Software-Defined Networking,'' SIGCOMM 2022.
\bibitem{titanium2024} Google Cloud Blog, ``Google Axion and Titanium: A New Era of Infrastructure Efficiency,'' 2024.
\bibitem{axion2024} Google Cloud, ``Google Axion Processors: Performance and Efficiency Redefined,'' 2024.
\bibitem{switch16} 邹银超（衍云），《通信那些事儿篇16——网络平台基石"交换机"》，阿里ATA，2022.
\bibitem{router17} 邹银超（衍云），《通信那些事儿篇17——网络平台基石"路由器"》，阿里ATA，2022.
\bibitem{chip18} 邹银超（衍云），《通信那些事儿篇18——交换机和路由器芯片》，阿里ATA，2022.
\bibitem{tna2021} 于诗，《Barefoot TNA架构及其在网关场景的应用》，阿里云Cstar，2021.
\bibitem{sdn2022} 王云毅（云遇），《路线之争："软件"SDN与"硬件"SDN》，阿里ATA，2022.
\bibitem{tswitch2014} bobqian（钱波），《网络开放浪潮下顺势而为的腾讯自研交换机》，腾讯KM，2014.
\bibitem{tswitch2020} wilsonwen（文权），《网络设备自研之路——自研交换机的今天、明天》，腾讯，2020.
\end{thebibliography}
